\documentclass[11pt,oneside]{article}
% Standalone ProveIt research note. The source intentionally has no external
% inputs so that the TeX/PDF pair is directly reproducible after extraction.
\usepackage[a4paper,margin=26mm,headheight=15pt]{geometry}
\usepackage[T1]{fontenc}
\usepackage[utf8]{inputenc}
\IfFileExists{libertinus.sty}{\usepackage{libertinus}}{\usepackage{lmodern}}
\usepackage{microtype}
\usepackage{amsmath,amssymb,amsthm,mathtools,mathrsfs,bm}
\usepackage{booktabs,longtable,tabularx,array}
\usepackage{graphicx}
\usepackage{xcolor}
\usepackage{enumitem}
\usepackage{fancyhdr}
\usepackage{xurl}
\usepackage{listings}
\usepackage[most]{tcolorbox}
\usepackage{hyperref}
\usepackage[nameinlink,noabbrev,capitalise]{cleveref}

\definecolor{linkblue}{RGB}{24,72,124}
\definecolor{deepgreen}{RGB}{23,102,69}
\definecolor{deepred}{RGB}{140,42,42}
\definecolor{deepgold}{RGB}{122,91,19}
\definecolor{shadegray}{RGB}{246,247,249}
\definecolor{rulegray}{RGB}{195,201,208}
\definecolor{palered}{RGB}{251,236,236}
\definecolor{palegold}{RGB}{251,245,230}
\definecolor{palegreen}{RGB}{233,245,243}
\definecolor{paleblue}{RGB}{236,243,251}

\hypersetup{
  hypertexnames=false,
  colorlinks=true,
  linkcolor=linkblue,
  citecolor=deepgreen,
  urlcolor=linkblue,
  pdftitle={Asymptotic Inversion of the Kinkelin-Bendersky K-Function},
  pdfauthor={Vladimir Reshetnikov},
  pdfsubject={Lambert-anchored transseries for the inverse generalized hyperfactorial and a general theory of x^p log x reversion},
  pdfkeywords={K-function, hyperfactorial, Kinkelin, Bendersky, Lambert W, inverse asymptotics, transseries, Lagrange-Burmann inversion, Barnes G-function}
}

\pagestyle{fancy}
\fancyhf{}
\fancyhead[L]{\small Inverse K-function transseries}
\fancyhead[R]{\small\nouppercase{\leftmark}}
\fancyfoot[C]{\thepage}
\renewcommand{\headrulewidth}{0.4pt}

\setcounter{tocdepth}{2}
\setcounter{secnumdepth}{3}
\setlist[itemize]{leftmargin=1.5em,itemsep=0.25ex,topsep=0.6ex}
\setlist[enumerate]{leftmargin=1.9em,itemsep=0.35ex,topsep=0.6ex}
\allowdisplaybreaks[2]

\theoremstyle{plain}
\newtheorem{theorem}{Theorem}[section]
\newtheorem{proposition}[theorem]{Proposition}
\newtheorem{lemma}[theorem]{Lemma}
\newtheorem{corollary}[theorem]{Corollary}
\theoremstyle{definition}
\newtheorem{definition}[theorem]{Definition}
\newtheorem{algorithm}[theorem]{Algorithm}
\newtheorem{example}[theorem]{Example}
\theoremstyle{remark}
\newtheorem{remark}[theorem]{Remark}

\crefname{algorithm}{algorithm}{algorithms}
\Crefname{algorithm}{Algorithm}{Algorithms}

\newtcolorbox{mainbox}[1][]{%
  enhanced,breakable,colback=paleblue,colframe=linkblue,
  boxrule=0.7pt,arc=1.2mm,left=2mm,right=2mm,top=1.4mm,bottom=1.4mm,
  title={Main result},fonttitle=\bfseries,#1}
\newtcolorbox{warningbox}[1][]{%
  enhanced,breakable,colback=palered,colframe=deepred,
  boxrule=0.6pt,arc=1.2mm,left=2mm,right=2mm,top=1.2mm,bottom=1.2mm,
  title={Scale and branch warning},fonttitle=\bfseries,#1}
\newtcolorbox{methodbox}[1][]{%
  enhanced,breakable,colback=palegreen,colframe=deepgreen,
  boxrule=0.6pt,arc=1.2mm,left=2mm,right=2mm,top=1.2mm,bottom=1.2mm,
  title={Reusable method},fonttitle=\bfseries,#1}
\newtcolorbox{summarybox}[1][]{%
  enhanced,breakable,colback=palegold,colframe=deepgold,
  boxrule=0.6pt,arc=1.2mm,left=2mm,right=2mm,top=1.2mm,bottom=1.2mm,
  title={Summary},fonttitle=\bfseries,#1}

\lstdefinestyle{pythonstyle}{%
  language=Python,
  basicstyle=\ttfamily\small,
  columns=fullflexible,
  frame=single,
  rulecolor=\color{rulegray},
  backgroundcolor=\color{shadegray},
  showstringspaces=false,
  keepspaces=true,
  breaklines=true,
  xleftmargin=0.4em,xrightmargin=0.4em,
  aboveskip=0.8em,belowskip=0.8em
}
\lstset{style=pythonstyle}

\newcommand{\W}{\operatorname{W}}
\newcommand{\e}{\mathrm{e}}
\newcommand{\dd}{\mathop{}\!\mathrm{d}}
\newcommand{\R}{\mathbb{R}}
\newcommand{\C}{\mathbb{C}}
\newcommand{\Q}{\mathbb{Q}}
\newcommand{\Acal}{\mathcal{A}}
\newcommand{\Fcal}{\mathcal{F}}
\newcommand{\Hcal}{\mathcal{H}}
\newcommand{\Lcal}{\mathcal{L}}
\newcommand{\Tcal}{\mathcal{T}}
\newcommand{\Bcal}{\mathcal{B}}
\newcommand{\coeff}[2]{\left[#1^{#2}\right]}
\DeclareMathOperator{\Res}{Res}
\DeclareMathOperator{\Log}{Log}
\DeclareMathOperator{\Trunc}{Trunc}

\begin{document}

\title{\textbf{Asymptotic Inversion of the Kinkelin--Bendersky \(K\)-Function}\\[0.35em]
  \Large Lambert-Anchored Transseries for the Generalized Hyperfactorial\\
  and a General Theory of \(x^p\log x\) Reversion}
\author{Vladimir Reshetnikov\\
  \small ProveIt research note}
\date{3 September 2026}
\maketitle

\begin{abstract}
The Kinkelin--Bendersky function \(K\), characterized on the positive real axis by
\(K(x+1)=x^xK(x)\) and \(K(1)=1\), interpolates the hyperfactorials
\(K(n+1)=\prod_{k=1}^n k^k\).  This article derives an all-orders asymptotic
transseries for its large real inverse and for the corresponding complex
sectorial inverse germs.  The decisive normalization is the half-shift
\(t=z-\tfrac12\).  With \(A\) the Glaisher--Kinkelin constant and
\(C=\log A-\tfrac18\), one has
\[
 \log K\!\left(t+\tfrac12\right)
 \sim \frac12t^2\log t-\frac14t^2+C-\frac1{24}\log t
      +\sum_{m\ge1}\frac{\kappa_m}{t^{2m}},
\]
where
\[
 \kappa_m=-\frac{B_{2m+2}(1/2)}{4m(m+1)(2m+1)}
 =\frac{(1-2^{-2m-1})B_{2m+2}}{4m(m+1)(2m+1)}.
\]
For \(y\to+\infty\), put \(Y=\log y\),
\[
 w=\W_0(4Y/\e),\qquad \lambda=\frac{w+1}{2},\qquad
 T=\e^\lambda=2\sqrt{\frac{Y}{w}},\qquad q=T^{-2}.
\]
Then the inverse branch mapping \([1,\infty)\) to \([2,\infty)\) satisfies
\[
 K_+^{-1}(y)=\frac12+T+\sum_{n=1}^{N}a_n(\lambda)T^{1-2n}
              +O_N(T^{-2N-1}),
\]
for every fixed \(N\).  Every \(a_n\) is an explicitly computable polynomial in
\(\lambda^{-1}\) of degree at most \(2n-1\), with coefficients in
\(\Q[C,\kappa_1,\ldots,\kappa_{n-1}]\).  We give a triangular recurrence, a
closed ordered-operator formula obtained from a perturbative
Lagrange--B\"urmann identity, four explicit correction blocks, a rank-two
reorganization in \((q,\lambda^{-1})\), rigorous residual-to-error transfer,
and high-precision numerical checks.

A second, accelerated Lambert anchor exactly absorbs the constant and logarithmic
corrections.  It gives the one-line approximation
\[
 K_+^{-1}(y)\approx \frac12+
 \sqrt{\frac{4(Y-C)}{\W_0(4(Y-C)/\e)}+\frac1{12}},
\]
whose inverse error is already \(O(T^{-3}/\log T)\).  Finally, the calculation is
placed in a general framework for maps
\(F(x)=a x^p(\log x-b)+H(x)\), where \(H\) is a lower-order grid-based
logarithmic transseries.  The Lambert anchor is exact, the inverse support is
locally finite, and the coefficients follow from a universal differential
operator calculus.
\end{abstract}

\noindent\textbf{Keywords.}
Kinkelin--Bendersky \(K\)-function; hyperfactorial; generalized hyperfactorial;
Glaisher--Kinkelin constant; Barnes \(G\)-function; Lambert \(W\); asymptotic
inverse; transseries; Lagrange--B\"urmann inversion; series reversion.

\begingroup
\footnotesize
\let\originaladdvspace\addvspace
\renewcommand{\addvspace}[1]{\originaladdvspace{0.25em}}
\tableofcontents
\endgroup
\clearpage

\section{The inverse problem and the answer in compact form}
\label{sec:intro}

The object of study is the real-valued Kinkelin--Bendersky function
\(K:[0,\infty)\to(0,\infty)\), normalized by
\begin{equation}
 K(0)=K(1)=K(2)=1,
 \qquad K(x+1)=x^xK(x)\quad(x>0).
 \label{eq:K-functional}
\end{equation}
At positive integers,
\begin{equation}
 K(n+1)=\prod_{k=1}^{n}k^k.
 \label{eq:K-integers}
\end{equation}
This is the convention used in the standard literature on the \(K\)-function
\cite{Bendersky1933,SondowHadjicostas2007}.  Some software and combinatorial
sources instead denote by \(H(x)\) the shifted hyperfactorial
\begin{equation}
 H(x):=K(x+1).
 \label{eq:H-shift}
\end{equation}
Thus an inverse formula for \(H\) is obtained from the formulas below by
subtracting \(1\) from the inverse of \(K\).

For large positive arguments, the relevant inverse is the increasing branch
\begin{equation}
 K_+^{-1}:[1,\infty)\longrightarrow[2,\infty).
 \label{eq:real-inverse-branch}
\end{equation}
The endpoint value \(y=1\) is not globally one-to-one because
\(K(0)=K(1)=K(2)=1\); the subscript \(+\) fixes the branch which grows to
infinity.

\begin{mainbox}
Let
\begin{equation}
 Y=\log y,
 \qquad w=\W_0\!\left(\frac{4Y}{\e}\right),
 \qquad \lambda=\frac{w+1}{2},
 \qquad T=\e^\lambda=2\sqrt{\frac{Y}{w}}.
 \label{eq:main-variables}
\end{equation}
Let \(A\) be the Glaisher--Kinkelin constant and set
\begin{equation}
 C=\log A-\frac18=-\zeta'(-1)-\frac1{24}.
 \label{eq:C-def}
\end{equation}
Then, as \(y\to+\infty\),
\begin{equation}
 K_+^{-1}(y)
 \sim \frac12+T+\sum_{n\ge1}a_n(\lambda)T^{1-2n}.
 \label{eq:main-transseries}
\end{equation}
The first three coefficients are
\begin{align}
 a_1(\lambda)={}&\frac1{24}-\frac{C}{\lambda},
 \label{eq:a1}\\[0.3em]
 a_2(\lambda)={}&-\frac1{1152}
 +\frac{20C+1}{480\lambda}
 -\frac{C^2}{2\lambda^2}
 -\frac{C^2}{2\lambda^3},
 \label{eq:a2}\\[0.3em]
 a_3(\lambda)={}&\frac1{27648}
 -\frac{1890C+311}{725760\lambda}
 +\frac{C(10C+1)}{160\lambda^2}
 \notag\\
 &-\frac{C(240C^2-10C-1)}{480\lambda^3}
 -\frac{5C^3}{6\lambda^4}
 -\frac{C^3}{2\lambda^5}.
 \label{eq:a3}
\end{align}
For every fixed \(N\ge0\), truncation after \(a_N\) has the real-axis error
\begin{equation}
 K_+^{-1}(y)-\left(\frac12+T+
       \sum_{n=1}^{N}a_n(\lambda)T^{1-2n}\right)
 =O_N(T^{-2N-1}).
 \label{eq:main-error}
\end{equation}
The convention is that the sum is empty when \(N=0\).
\end{mainbox}

The first correction can be written directly in \(Y\) and \(w\):
\begin{equation}
 K_+^{-1}(y)
 =\frac12+2\sqrt{\frac{Y}{w}}
 +\frac12\sqrt{\frac{w}{Y}}
   \left(\frac1{24}-\frac{2C}{w+1}\right)
 +O\!\left(\left(\frac{w}{Y}\right)^{3/2}\right).
 \label{eq:first-correction-Yw}
\end{equation}
This already displays the two essential scales.  The Lambert variable
\(w\sim\log Y\) resums the entire iterated-logarithm expansion of the leading
block.  The first genuinely new block is smaller by
\begin{equation}
 T^{-2}=\e^{-(w+1)}=\frac{w}{4Y},
 \label{eq:q-def-intro}
\end{equation}
which is beyond every algebraic order in \(1/w\).

\subsection{What is meant by a transseries here}

There are three equivalent ways to read \cref{eq:main-transseries}.

\begin{enumerate}
\item As a \emph{Lambert-anchored outer expansion}, it is a Poincar\'e expansion
in powers of \(q=T^{-2}\), with coefficient functions of the slow variable
\(\lambda=\log T\).
\item As a \emph{rank-two grid-based transseries}, its monomials are
\begin{equation}
 Tq^n\lambda^{-j}
 =\exp\!\left(-\frac{2n-1}{2}(w+1)\right)\lambda^{-j},
 \qquad n\ge1,
 \label{eq:rank-two-monomials}
\end{equation}
with \(0\le j\le2n-1\) in the present problem.
\item After expanding \(w=\W_0(4Y/\e)\) in
\(L_1=\log(4Y/\e)\) and \(L_2=\log L_1\), it becomes a nested
logarithmic transseries in powers of \(L_1/Y\), whose coefficient blocks are
polynomials in \(L_2\) and inverse powers of \(L_1\).
\end{enumerate}

The Lambert-anchored form is both shorter and numerically superior.  Expanding
\(\W\) too early destroys an exact inversion and replaces it with an infinite
slow series before the first algebraic correction has even appeared.

\subsection{Outline and status of the results}

The forward asymptotic expansion of the hyperfactorial and of related
Bendersky functions is classical and has been refined in a substantial
literature; see, for example,
\cite{Bendersky1933,WangLiu2016,Wang2017,Xu2020}.  A useful closed Lambert
approximation to the inverse has also appeared in an online discussion
\cite{MathSEHyperfactorial}.  The all-orders inverse coefficients, the centered
operator formula, the rank-two resummation, and the general
\(a x^p(\log x-b)\) inversion calculus are derived here.  This statement is a
description of the contents of this note, not a claim that no equivalent
formula can exist elsewhere.

The argument deliberately separates three levels of validity.

\begin{itemize}
\item The identities for coefficients are exact formal identities in a
locally finite transseries algebra.
\item On the positive real axis, the source expansion and monotonicity turn
each fixed truncation into the rigorous asymptotic estimate
\cref{eq:main-error}.
\item In the complex plane, the same formulas describe local inverse germs on
admissible sectors, once the logarithm and Lambert branches are fixed and the
path stays away from critical points.
\end{itemize}

\section{The \(K\)-function and its increasing real branch}
\label{sec:K-background}

\subsection{Integral, zeta, and Barnes representations}

For \(x\ge0\), one convenient definition is
\begin{equation}
 \log K(x)
 =\frac{x^2-x}{2}-\frac{x}{2}\log(2\pi)
  +\int_0^x\log\Gamma(u)\,\dd u.
 \label{eq:K-integral}
\end{equation}
The singularity of \(\log\Gamma(u)\) at the origin is integrable.  A standard
complex continuation is encoded by
\begin{equation}
 K(z)=\exp\!\bigl(\zeta'(-1,z)-\zeta'(-1)\bigr),
 \label{eq:K-Hurwitz}
\end{equation}
with compatible branch conventions.  On the positive real axis, the relation
to the Barnes \(G\)-function is
\begin{equation}
 K(z)G(z)=\Gamma(z)^{z-1}.
 \label{eq:K-Barnes-relation}
\end{equation}
Equivalently,
\begin{equation}
 K(x+1)G(x+1)=\Gamma(x+1)^x.
 \label{eq:K-Barnes-shifted}
\end{equation}
These formulas and the recurrence \cref{eq:K-functional} are recorded in
\cite{SondowHadjicostas2007}; the Barnes asymptotics used below are summarized
in the NIST Digital Library of Mathematical Functions
\cite{DLMF}.

\subsection{Monotonicity beyond \(2\)}

Differentiate \cref{eq:K-integral}:
\begin{equation}
 \frac{\dd}{\dd x}\log K(x)
 =x-\frac12-\frac12\log(2\pi)+\log\Gamma(x),
 \label{eq:logK-prime}
\end{equation}
so that
\begin{equation}
 \frac{\dd^2}{\dd x^2}\log K(x)=1+\psi(x),
 \label{eq:logK-second}
\end{equation}
where \(\psi=\Gamma'/\Gamma\).

\begin{proposition}[Increasing real branch]
\label{prop:monotone}
The function \(K\) is strictly increasing on \([2,\infty)\), maps this interval
onto \([1,\infty)\), and therefore has the inverse branch
\cref{eq:real-inverse-branch}.
\end{proposition}

\begin{proof}
The digamma function is increasing on \((0,\infty)\), and
\(\psi(2)=1-\gamma>0\).  Hence \(1+\psi(x)>0\) for \(x\ge2\), so the right-hand
side of \cref{eq:logK-prime} is increasing there.  At \(x=2\), it equals
\begin{equation*}
 \frac32-\frac12\log(2\pi)>0.
\end{equation*}
Thus \((\log K)'(x)>0\) for \(x\ge2\), and therefore \(K'(x)>0\).  Since
\(K(2)=1\) and the large-\(x\) expansion in \cref{sec:centered-source} tends to
infinity, the range is \([1,\infty)\).
\end{proof}

\begin{corollary}[Shifted hyperfactorial inverse]
\label{cor:H-inverse}
If \(H(x)=K(x+1)\), then on the increasing branch
\begin{equation}
 H_+^{-1}(y)=K_+^{-1}(y)-1
 \sim -\frac12+T+\sum_{n\ge1}a_n(\lambda)T^{1-2n}.
 \label{eq:H-inverse}
\end{equation}
\end{corollary}

\section{The centered forward expansion}
\label{sec:centered-source}

The forward asymptotic expansion is not merely an input to inversion.  Its
correct centering determines the support and the complexity of the inverse
transseries.

\subsection{The classical uncentered form}

The Barnes expansion
\cite[\S5.17(ii)]{DLMF} and the exact cancellation in
\cref{eq:K-Barnes-shifted} give, uniformly in every closed sector
\(|\arg z|\le\pi-\delta\),
\begin{align}
 \log K(z+1)
 \sim{}&\left(\frac{z^2}{2}+\frac z2+\frac1{12}\right)\log z
       -\frac{z^2}{4}+\log A
 \notag\\
 &-\sum_{m\ge1}
 \frac{B_{2m+2}}{4m(m+1)(2m+1)z^{2m}}.
 \label{eq:uncentered-logK}
\end{align}
Here \(B_n\) denotes the Bernoulli number and
\begin{equation}
 \log A=\frac1{12}-\zeta'(-1).
 \label{eq:logA}
\end{equation}
Although \cref{eq:uncentered-logK} contains only even negative powers, its
leading polynomial-logarithmic part contains a linear term \(z\log z\).  The
inverse becomes substantially simpler after the half-shift.

\subsection{Half-shift and parity}

Define
\begin{equation}
 t=z-\frac12,
 \qquad
 \Lcal(t):=\log K\!\left(t+\frac12\right),
 \qquad
 \Phi(t):=\frac12t^2\log t-\frac14t^2.
 \label{eq:centered-defs}
\end{equation}
Let the Bernoulli polynomials be normalized by
\begin{equation}
 \frac{u\e^{xu}}{\e^u-1}
 =\sum_{n\ge0}B_n(x)\frac{u^n}{n!}.
 \label{eq:Bernoulli-polys}
\end{equation}
They satisfy
\begin{equation}
 B_n\!\left(\frac12\right)=(2^{1-n}-1)B_n.
 \label{eq:Bernoulli-half}
\end{equation}

\begin{theorem}[Canonical centered expansion]
\label{thm:centered-expansion}
As \(t\to\infty\) in every closed sector
\(|\arg t|\le\pi-\delta\),
\begin{equation}
 \Lcal(t)
 \sim \Phi(t)+C-\frac1{24}\log t
       +\sum_{m\ge1}\kappa_m t^{-2m},
 \label{eq:centered-logK}
\end{equation}
where \(C\) is given by \cref{eq:C-def} and
\begin{equation}
 \kappa_m
 =-\frac{B_{2m+2}(1/2)}{4m(m+1)(2m+1)}
 =\frac{(1-2^{-2m-1})B_{2m+2}}
        {4m(m+1)(2m+1)}.
 \label{eq:kappa-formula}
\end{equation}
For every fixed \(M\ge1\), the corresponding finite form is
\begin{equation}
 \Lcal(t)=\Phi(t)+C-\frac1{24}\log t
 +\sum_{m=1}^{M-1}\kappa_m t^{-2m}+O(t^{-2M}).
 \label{eq:centered-remainder}
\end{equation}
\end{theorem}

\begin{proof}
From \cref{eq:logK-prime}, with \(x=t+\tfrac12\),
\begin{equation}
 \Lcal'(t)=t-\frac12\log(2\pi)+\log\Gamma\!\left(t+\frac12\right).
 \label{eq:Lprime-gamma}
\end{equation}
The shifted Stirling expansion is
\begin{equation}
 \log\Gamma(t+a)
 \sim \left(t+a-\frac12\right)\log t-t+\frac12\log(2\pi)
 +\sum_{n\ge1}
   \frac{(-1)^{n+1}B_{n+1}(a)}{n(n+1)t^n}.
 \label{eq:shifted-Stirling}
\end{equation}
At \(a=\tfrac12\), all Bernoulli polynomials of odd degree greater than one
vanish.  Hence
\begin{equation}
 \Lcal'(t)
 \sim t\log t-\frac1{24t}
 +\sum_{m\ge1}
 \frac{B_{2m+2}(1/2)}{(2m+1)(2m+2)t^{2m+1}}.
 \label{eq:centered-derivative-expansion}
\end{equation}
Termwise integration gives \cref{eq:centered-logK} up to a constant and yields
exactly \cref{eq:kappa-formula}.  Matching with the uncentered expansion
\cref{eq:uncentered-logK} determines the constant as
\(\log A-\tfrac18=C\).  Uniform remainder estimates follow either by applying
the standard sectorial remainder for shifted Stirling and integrating, or
directly by re-expanding the Barnes remainder.
\end{proof}

The first coefficients are
\begin{equation}
 \kappa_1=-\frac7{5760},\quad
 \kappa_2=\frac{31}{161280},\quad
 \kappa_3=-\frac{127}{1290240},\quad
 \kappa_4=\frac{511}{4866048},\quad
 \kappa_5=-\frac{1414477}{7380172800}.
 \label{eq:kappa-first-five}
\end{equation}

\begin{remark}[Why the half-shift is canonical]
The half-shift simultaneously removes the \(z\log z\) term from the dominant
phase and converts the entire tail to even inverse powers.  At the derivative
level this is the identity \(B_{2m+1}(1/2)=0\).  Consequently the inverse
corrections occur in the sparse lattice \(T^{1-2n}\), rather than in every
integer power of \(T^{-1}\).
\end{remark}

\begin{warningbox}
The variable in \cref{eq:centered-logK} is
\(t=z-\tfrac12\), where \(z\) is the argument of \(K\).  If one instead uses
the common hyperfactorial variable \(x\) in \(H(x)=K(x+1)\), then
\(t=x+\tfrac12\).  Losing this shift produces an apparent error of order
\(1/2\), much larger than every displayed correction block.
\end{warningbox}

\section{Exact Lambert inversion of the dominant phase}
\label{sec:Lambert-anchor}

The leading phase is
\begin{equation}
 \Phi(t)=\frac14t^2\bigl(\log(t^2)-1\bigr).
 \label{eq:Phi-alt}
\end{equation}
It is not necessary to invert this phase asymptotically: it is exactly
Lambert-invertible.

\begin{proposition}[Exact dominant anchor]
\label{prop:exact-anchor}
Let \(Y>0\) and define \(w,\lambda,T\) by \cref{eq:main-variables}.  Then
\begin{equation}
 \Phi(T)=Y,
 \qquad \log T=\lambda,
 \qquad T^2=\frac{4Y}{w},
 \qquad T^{-2}=\e^{-(w+1)}=\frac{w}{4Y}.
 \label{eq:anchor-identities}
\end{equation}
\end{proposition}

\begin{proof}
Set \(w=2\log T-1\).  Then
\begin{equation}
 4\Phi(T)=T^2(2\log T-1)=\e^{w+1}w.
\end{equation}
Thus \(\Phi(T)=Y\) is equivalent to
\begin{equation}
 w\e^w=\frac{4Y}{\e}.
\end{equation}
On the positive real axis this gives \(w=\W_0(4Y/\e)\), followed by all the
identities in \cref{eq:anchor-identities}.
\end{proof}

\begin{corollary}[Leading inverse]
\label{cor:leading-inverse}
As \(y\to+\infty\),
\begin{equation}
 K_+^{-1}(y)=\frac12+T+O(T^{-1}).
 \label{eq:leading-inverse}
\end{equation}
\end{corollary}

\begin{proof}
The omitted part of \cref{eq:centered-logK} is \(O(\log T)\), while
\(\Phi'(T)=T\log T\).  One inverse-function step therefore produces a
displacement \(O(T^{-1})\).
\end{proof}

The two slow-fast parameters used throughout the rest of the article are
\begin{equation}
 \lambda=\log T=\frac{w+1}{2},
 \qquad q=T^{-2}=\e^{-2\lambda}.
 \label{eq:lambda-q}
\end{equation}
They satisfy the strict scale separation
\begin{equation}
 q=o(\lambda^{-M})\qquad(M=0,1,2,\ldots).
 \label{eq:q-beyond-all-orders}
\end{equation}
Thus a correction of relative size \(q\) is beyond all orders of the ordinary
large-\(\lambda\) expansion.

\section{The normalized inverse equation and triangular reversion}
\label{sec:implicit-equation}

Write the unknown centered inverse in the form
\begin{equation}
 t=T(1+u),
 \qquad
 u\sim\sum_{n\ge1}a_n(\lambda)q^n,
 \qquad q=T^{-2}.
 \label{eq:u-ansatz}
\end{equation}
Because \(\Phi(T)=Y\), direct substitution gives
\begin{equation}
 \frac{\Phi(T(1+u))-\Phi(T)}{T^2}
 =\Psi_\lambda(u),
 \label{eq:Psi-origin}
\end{equation}
where
\begin{equation}
 \Psi_\lambda(u)
 :=\frac12(1+u)^2
       \left(\lambda+\log(1+u)-\frac12\right)
   -\frac12\left(\lambda-\frac12\right).
 \label{eq:Psi-def}
\end{equation}
In particular,
\begin{equation}
 \Psi_\lambda(0)=0,
 \qquad \partial_u\Psi_\lambda(0)=\lambda.
 \label{eq:Psi-linear}
\end{equation}

Let
\begin{equation}
 \Hcal(q,\lambda)
 :=C-\frac{\lambda}{24}+\sum_{m\ge1}\kappa_mq^m.
 \label{eq:Hcal-K}
\end{equation}
The complete formal inverse equation is
\begin{align}
 0={}&\Fcal(u,q;\lambda)
 \notag\\
 :={}&\Psi_\lambda(u)
 +q\left(C-\frac{\lambda+\log(1+u)}{24}\right)
 +\sum_{m\ge1}\kappa_mq^{m+1}(1+u)^{-2m}.
 \label{eq:normalized-implicit}
\end{align}
The equality is to be read formally when the full divergent source series is
used, and as an asymptotic identity to any prescribed finite order when
\cref{eq:centered-remainder} is truncated.

\begin{theorem}[Unique formal inverse and triangular recurrence]
\label{thm:triangular-recurrence}
There is a unique formal solution
\begin{equation}
 U(q;\lambda)=\sum_{n\ge1}a_n(\lambda)q^n
 \label{eq:U-formal}
\end{equation}
to \(\Fcal(U,q;\lambda)=0\) with \(U(0;\lambda)=0\), after adjoining
\(\lambda^{-1}\) to the coefficient ring.  If
\begin{equation}
 U_{N-1}(q;\lambda)=\sum_{j=1}^{N-1}a_j(\lambda)q^j,
 \label{eq:UNminus1}
\end{equation}
then
\begin{equation}
 a_N(\lambda)
 =-\frac1\lambda\,[q^N]\,
   \Fcal\!\left(U_{N-1}(q;\lambda),q;\lambda\right).
 \label{eq:triangular-recurrence}
\end{equation}
Only \(\kappa_1,\ldots,\kappa_{N-1}\) are needed to compute \(a_N\).
\end{theorem}

\begin{proof}
Suppose \(a_1,\ldots,a_{N-1}\) have already been determined.  Replace
\(U_{N-1}\) by \(U_{N-1}+a_Nq^N\) in \cref{eq:normalized-implicit}.  At order
\(q^N\), the new unknown can enter only through the linearization of
\(\Psi_\lambda\) at \(u=0\); all nonlinear occurrences multiply it by an
additional positive power of \(q\).  By \cref{eq:Psi-linear}, its coefficient
is \(\lambda a_N\).  Setting the coefficient of \(q^N\) to zero gives
\cref{eq:triangular-recurrence}.  This also proves uniqueness by induction.
The source term \(\kappa_mq^{m+1}\) can affect order \(q^N\) only when
\(m\le N-1\).
\end{proof}

The recurrence is often the fastest way to generate coefficients in a computer
algebra system.  It also reveals why the expansion is locally finite: no
coefficient requires an infinite part of the forward asymptotic series.

\subsection{A polynomial structure theorem}

It is useful to isolate the dependence on the second small parameter
\begin{equation}
 s=\lambda^{-1}.
 \label{eq:s-def}
\end{equation}
Decompose
\begin{equation}
 \Psi_\lambda(u)=\lambda A(u)+B(u),
 \label{eq:Psi-AB}
\end{equation}
where
\begin{align}
 A(u)&=u+\frac{u^2}{2},
 \label{eq:A-def}\\
 B(u)&=\frac12(1+u)^2\log(1+u)-\frac u2-\frac{u^2}{4}.
 \label{eq:B-def}
\end{align}
After division by \(\lambda\), the implicit equation becomes
\begin{equation}
 A(u)-\frac q{24}+sR(u,q)=0,
 \label{eq:s-implicit}
\end{equation}
with
\begin{equation}
 R(u,q)=B(u)+q\left(C-\frac{\log(1+u)}{24}\right)
 +\sum_{m\ge1}\kappa_mq^{m+1}(1+u)^{-2m}.
 \label{eq:R-def}
\end{equation}
All explicit dependence on \(\lambda\) has now been compressed into one factor
\(s\).

\begin{theorem}[Finite slow depth of each outer block]
\label{thm:finite-slow-depth}
For every \(N\ge1\),
\begin{equation}
 a_N(\lambda)
 \in\Q[C,\kappa_1,\ldots,\kappa_{N-1}][\lambda^{-1}],
 \qquad
 \deg_{\lambda^{-1}}a_N\le2N-1.
 \label{eq:aN-ring}
\end{equation}
Moreover,
\begin{equation}
 [\lambda^0]a_N(\lambda)
 =\binom{1/2}{N}12^{-N}.
 \label{eq:aN-constant-term}
\end{equation}
\end{theorem}

\begin{proof}
At \(s=0\), \cref{eq:s-implicit} reduces to
\begin{equation}
 u+\frac{u^2}{2}=\frac q{24}.
\end{equation}
The solution with \(u(0)=0\) is
\begin{equation}
 u=\sqrt{1+\frac q{12}}-1,
 \label{eq:u-szero}
\end{equation}
which proves \cref{eq:aN-constant-term}.  For the degree bound, use induction
on \(N\) in the coefficient equation obtained from
\cref{eq:s-implicit}.  The term linear in \(a_N\) has coefficient one.  A
product of previously known coefficients with total outer degree \(N\) has
slow degree at most \(2N-2\); multiplication by the single explicit factor
\(s\) in \(sR\) raises this to at most \(2N-1\).  Formal composition with
\(\log(1+u)\) and \((1+u)^{-2m}\) preserves the same outer-degree estimate.
The coefficient field is generated only by the source constants that can occur
through order \(q^N\).
\end{proof}

\begin{remark}
The finiteness in \cref{eq:aN-ring} is special and stronger than what is needed
for a transseries.  A general coefficient differential algebra may produce an
infinite slow expansion inside each outer block.  Here the linear appearance of
\(\lambda\) in the centered source equation forces a polynomial in
\(\lambda^{-1}\).
\end{remark}

\section{Perturbative Lagrange--B\"urmann inversion}
\label{sec:Lagrange}

The triangular recurrence is constructive, but an all-orders closed formula is
more revealing.  The needed identity is a perturbative form of
Lagrange--B\"urmann inversion; related uses of Lagrange inversion in
asymptotic factorial problems appear in \cite{BrassescoMendez2011}.

\subsection{A general differential-operator identity}

\begin{theorem}[Perturbative Lagrange--B\"urmann formula]
\label{thm:perturbative-LB}
Let \(F_0\) and \(H\) be analytic near \(X\), suppose
\(F_0'(X)\ne0\), and let \(x(\varepsilon)\) be the local solution of
\begin{equation}
 F_0(x)+\varepsilon H(x)=F_0(X),
 \qquad x(0)=X.
 \label{eq:perturbed-equation}
\end{equation}
Define
\begin{equation}
 \mathscr D_0:=\frac1{F_0'(X)}\frac{\dd}{\dd X}.
 \label{eq:D0-general}
\end{equation}
Then, as a convergent local power series in \(\varepsilon\), and also as a
formal identity,
\begin{equation}
 x(\varepsilon)
 =X+\sum_{n\ge1}\frac{(-\varepsilon)^n}{n!}
 \mathscr D_0^{\,n-1}
 \left(\frac{H(X)^n}{F_0'(X)}\right).
 \label{eq:perturbative-LB}
\end{equation}
More generally, for analytic \(g\),
\begin{equation}
 g(x(\varepsilon))
 =g(X)+\sum_{n\ge1}\frac{(-\varepsilon)^n}{n!}
 \mathscr D_0^{\,n-1}
 \left(\frac{g'(X)H(X)^n}{F_0'(X)}\right).
 \label{eq:perturbative-LB-g}
\end{equation}
\end{theorem}

\begin{proof}
Take a small positively oriented contour around \(X\) containing no other zero
of \(F_0(z)-F_0(X)\).  For sufficiently small \(\varepsilon\), the argument
principle gives
\begin{equation}
 g(x(\varepsilon))
 =\frac{1}{2\pi i}\oint
 g(z)\frac{F_0'(z)+\varepsilon H'(z)}
 {F_0(z)+\varepsilon H(z)-F_0(X)}\,\dd z.
 \label{eq:argument-principle-g}
\end{equation}
Subtract the same formula at \(\varepsilon=0\), write the difference as a
logarithmic derivative, integrate by parts, and expand
\(\log(1+\varepsilon H/(F_0-F_0(X)))\).  The coefficient of
\(\varepsilon^n\) is
\begin{equation}
 \frac{(-1)^n}{n}
 \Res_{z=X}\frac{g'(z)H(z)^n}{(F_0(z)-F_0(X))^n}.
 \label{eq:residue-step}
\end{equation}
The change of local coordinate \(v=F_0(z)\) yields the standard residue
identity
\begin{equation}
 \Res_{z=X}\frac{\varphi(z)}{(F_0(z)-F_0(X))^n}
 =\frac1{(n-1)!}
 \left.
 \left(\frac1{F_0'(z)}\frac{\dd}{\dd z}\right)^{n-1}
 \frac{\varphi(z)}{F_0'(z)}
 \right|_{z=X}.
 \label{eq:residue-coordinate}
\end{equation}
Combining the two formulas proves \cref{eq:perturbative-LB-g}; taking
\(g(z)=z\) gives \cref{eq:perturbative-LB}.
\end{proof}

\begin{methodbox}
To invert \(F=F_0+H\), first solve \(F_0(X)=Y\) exactly or asymptotically, then
apply \cref{eq:perturbative-LB} with \(\varepsilon=1\).  If the operator
\(\mathscr D_0\) strictly raises a chosen asymptotic valuation, the infinite
sum is locally finite coefficient by coefficient and is valid as a formal
transseries even when it is not a convergent perturbation series.
\end{methodbox}

\subsection{Specialization to the \(K\)-function}

For the centered \(K\)-problem,
\begin{equation}
 F_0(T)=\Phi(T),
 \qquad
 F_0'(T)=T\lambda,
 \qquad
 H(T)=C-\frac{\lambda}{24}+\sum_{m\ge1}\kappa_mT^{-2m}.
 \label{eq:K-F0-H}
\end{equation}
Therefore
\begin{equation}
 \mathscr D:=\frac1{T\lambda}\frac{\dd}{\dd T}
 \label{eq:K-D}
\end{equation}
and the centered inverse has the formal expansion
\begin{equation}
 t
 \sim T+\sum_{n\ge1}\frac{(-1)^n}{n!}
 \mathscr D^{\,n-1}
 \left(\frac{H(T)^n}{T\lambda}\right).
 \label{eq:K-LB-master}
\end{equation}
Equation \cref{eq:K-LB-master} is an all-orders formula requiring no implicit
coefficient solving.

To make the outer blocks explicit, observe that
\begin{equation}
 \mathscr D\bigl(Tq^r f(\lambda)\bigr)
 =Tq^{r+1}\Acal_r f(\lambda),
 \label{eq:D-block-action}
\end{equation}
where
\begin{equation}
 \Acal_r f
 :=\frac{f'(\lambda)+(1-2r)f(\lambda)}{\lambda}.
 \label{eq:Ar-def}
\end{equation}
Here and below, a prime on a coefficient function means differentiation with
respect to \(\lambda\).

\begin{theorem}[Ordered block formula]
\label{thm:block-formula}
Let
\begin{equation}
 \Hcal(q,\lambda)=C-\frac{\lambda}{24}
                   +\sum_{m\ge1}\kappa_mq^m.
 \label{eq:Hcal-repeated}
\end{equation}
Then the coefficient of \(Tq^N=T^{1-2N}\) in the inverse is
\begin{align}
 a_N(\lambda)
 ={}&\sum_{n=1}^{N}\frac{(-1)^n}{n!}
 \bigl(\Acal_{N-1}\circ\Acal_{N-2}\circ\cdots
        \circ\Acal_{N-n+1}\bigr)
 \notag\\
 &\hspace{28mm}\times
 \left(
   \frac{[q^{N-n}]\Hcal(q,\lambda)^n}{\lambda}
 \right).
 \label{eq:block-formula}
\end{align}
For \(n=1\), the displayed operator composition is empty and acts as the
identity.  In every nonempty composition, the rightmost operator acts first.
\end{theorem}

\begin{proof}
In the \(n\)-th summand of \cref{eq:K-LB-master}, a term
\([q^m]\Hcal^n\) begins as \(Tq^{m+1}\).  Each application of
\(\mathscr D\) raises the outer index by one according to
\cref{eq:D-block-action}.  After \(n-1\) applications, the final index is
\(m+n\).  Setting \(m=N-n\) and recording the ordered operators gives
\cref{eq:block-formula}.  Since \(n\le N\), only finitely many perturbation
orders contribute to a fixed block.
\end{proof}

\subsection{The first four inverse blocks}

For reference, the next coefficient after \cref{eq:a1}--\cref{eq:a3} is
\begin{align}
 a_4(\lambda)={}&-\frac5{2654208}
 +\frac{6300C+5629}{34836480\lambda}
 -\frac{37800C^2+12440C+63}{5806080\lambda^2}
 \notag\\
 &+\frac{3024000C^3+415800C^2-12440C-63}
        {29030400\lambda^3}
 \notag\\
 &-\frac{C^2(1800C^2-220C-33)}{2880\lambda^4}
 -\frac{C^2(1360C^2-20C-3)}{960\lambda^5}
 \notag\\
 &-\frac{35C^4}{24\lambda^6}
 -\frac{5C^4}{8\lambda^7}.
 \label{eq:a4}
\end{align}
The constant terms
\begin{equation}
 \frac1{24},\quad -\frac1{1152},\quad
 \frac1{27648},\quad -\frac5{2654208},\ldots
 \label{eq:constant-blocks}
\end{equation}
are precisely the coefficients of
\(\sqrt{1+q/12}-1\), as predicted by
\cref{eq:aN-constant-term}.

A compact alternative notation is
\begin{equation}
 \eta:=\frac{\lambda}{24}-C.
 \label{eq:eta-def}
\end{equation}
Then
\begin{align}
 a_1={}&\frac{\eta}{\lambda},
 \label{eq:a1-eta}\\
 a_2={}&-\frac{\eta^2}{2\lambda^3}
 -\frac{\eta(12\eta-1)}{24\lambda^2}
 +\frac7{5760\lambda},
 \label{eq:a2-eta}\\
 a_3={}&\frac{\eta^3}{2\lambda^5}
 +\frac{\eta^2(40\eta-3)}{48\lambda^4}
 +\frac{\eta(960\eta^2-160\eta+1)}{1920\lambda^3}
 \notag\\
 &-\frac{7(72\eta-1)}{138240\lambda^2}
 -\frac{31}{161280\lambda}.
 \label{eq:a3-eta}
\end{align}
This version makes the direct contributions of \(\kappa_1\) and
\(\kappa_2\) more visible.

\section{Analytic validity, remainder transfer, and certification}
\label{sec:analytic-validity}

Formal reversion becomes an asymptotic theorem because the source expansion is
uniform and the derivative of the dominant phase is large.

\subsection{Residual of a formal truncation}

For \(N\ge0\), define
\begin{equation}
 t_N:=T+\sum_{n=1}^{N}a_n(\lambda)T^{1-2n},
 \qquad
 z_N:=\frac12+t_N.
 \label{eq:tN-zN}
\end{equation}

\begin{lemma}[Residual order]
\label{lem:residual-order}
For every fixed \(N\ge0\),
\begin{equation}
 \log K(z_N)-Y=O_N(\lambda T^{-2N})
 \qquad(Y\to+\infty).
 \label{eq:residual-order}
\end{equation}
\end{lemma}

\begin{proof}
Divide the source equation by \(T^2\) as in
\cref{eq:normalized-implicit}.  By construction, the coefficients of
\(q,q^2,\ldots,q^N\) vanish after substituting the truncated series for \(u\).
The first possible uncancelled normalized term is of order
\(\lambda q^{N+1}\); the factor \(\lambda\) is present because the linear
part of \(\Psi_\lambda\) is \(\lambda u\).  Multiplication by
\(T^2=q^{-1}\) gives \(O(\lambda T^{-2N})\).  The source remainder
\(O(t^{-2N})\) in \cref{eq:centered-remainder} has the same or smaller order.
\end{proof}

\subsection{Transfer from residual to inverse error}

The centered derivative satisfies
\begin{equation}
 \Lcal'(t)=t\log t+O(t^{-1})
 \label{eq:Lprime-leading}
\end{equation}
uniformly on the positive real axis.  Since \(t_N/T\to1\),
\begin{equation}
 \Lcal'(t)=T\lambda\bigl(1+o(1)\bigr)
 \label{eq:Lprime-near-T}
\end{equation}
uniformly between \(t_N\) and the exact root for all sufficiently large
\(Y\).

\begin{theorem}[Real-axis asymptotic inverse]
\label{thm:real-asymptotic}
Let \(t_*:=K_+^{-1}(y)-\tfrac12\).  For every fixed \(N\ge0\),
\begin{equation}
 t_*-t_N=O_N(T^{-2N-1}).
 \label{eq:t-error}
\end{equation}
Equivalently, \cref{eq:main-error} holds.
\end{theorem}

\begin{proof}
By the mean value theorem, for some point \(\xi\) between \(t_N\) and \(t_*\),
\begin{equation}
 t_N-t_*
 =\frac{\Lcal(t_N)-\Lcal(t_*)}{\Lcal'(\xi)}
 =\frac{\log K(z_N)-Y}{\Lcal'(\xi)}.
 \label{eq:MVT-transfer}
\end{equation}
Use \cref{eq:residual-order,eq:Lprime-near-T} to obtain
\(O(\lambda T^{-2N})/(T\lambda)=O(T^{-2N-1})\).
\end{proof}

\begin{remark}[How many forward coefficients are required?]
To construct \(a_1,\ldots,a_N\), it is enough to know the centered forward
series through \(\kappa_{N-1}t^{-2N+2}\).  The next source term contributes at
the same scale as the first omitted inverse block and therefore belongs in the
remainder.
\end{remark}

\subsection{A posteriori residual certificates}

Asymptotic order does not by itself certify a finite numerical computation.
Monotonicity supplies a direct certificate.

\begin{proposition}[Residual-to-error bound]
\label{prop:residual-certificate}
Let \(Y>0\), let \(z_*=K_+^{-1}(\e^Y)\), and suppose an interval
\(I=[a,b]\subset[2,\infty)\) contains both \(z_*\) and an approximation
\(\widehat z\).  Then
\begin{equation}
 |\widehat z-z_*|
 \le
 \frac{|\log K(\widehat z)-Y|}
      {\inf_{x\in I}\left(x-\frac12-\frac12\log(2\pi)+\log\Gamma(x)\right)}.
 \label{eq:residual-certificate}
\end{equation}
The denominator is positive.  Since the derivative is increasing on
\([2,\infty)\), it may be replaced by its value at \(a\).
\end{proposition}

\begin{proof}
Apply the mean value theorem to \(\log K\) and use
\cref{eq:logK-prime,prop:monotone}.
\end{proof}

A first-order residual correction is
\begin{equation}
 z_*=\widehat z-
 \frac{\log K(\widehat z)-Y}
 {\widehat z-\frac12-\frac12\log(2\pi)+\log\Gamma(\widehat z)}
 +O\!\left(
   \frac{R^2\sup_I|1+\psi|}{(\inf_I(\log K)')^3}
 \right),
 \label{eq:Newton-residual}
\end{equation}
where \(R=\log K(\widehat z)-Y\).  Thus the transseries is also an excellent
starting value for Newton or Halley iteration.

\begin{warningbox}
All asymptotic series in this article are Poincar\'e expansions.  The Bernoulli
coefficients grow factorially, so summing indefinitely at fixed \(y\) is not
justified.  For high numerical accuracy, truncate near the least term or use a
small fixed number of inverse blocks followed by a certified Newton step.
\end{warningbox}

\section{Rank-two structure and logarithmic re-expansion}
\label{sec:rank-two}

The Lambert anchor packages an infinite logarithmic expansion into the single
variable \(w\).  This section unpacks the structure and also gives a second
resummation that reverses the order of the two scales.

\subsection{The ordinary large-argument expansion of \(\W\)}

Let
\begin{equation}
 L_1=\log\!\left(\frac{4Y}{\e}\right),
 \qquad L_2=\log L_1.
 \label{eq:L1-L2}
\end{equation}
The principal Lambert function has the classical expansion
\cite{CorlessEtAl1996,DLMF}
\begin{align}
 w={}&L_1-L_2+\frac{L_2}{L_1}
 +\frac{L_2(-2+L_2)}{2L_1^2}
 \notag\\
 &+\frac{L_2(6-9L_2+2L_2^2)}{6L_1^3}
 +O\!\left(\frac{L_2^4}{L_1^4}\right).
 \label{eq:W-log-expansion}
\end{align}
Consequently,
\begin{align}
 T=2\sqrt{\frac{Y}{L_1}}\biggl[{}&1+\frac{L_2}{2L_1}
 +\frac{L_2(3L_2-4)}{8L_1^2}
 \notag\\
 &+\frac{L_2(5L_2^2-16L_2+8)}{16L_1^3}
 \notag\\
 &+\frac{L_2(105L_2^3-568L_2^2+720L_2-192)}{384L_1^4}
 +O\!\left(\frac{L_2^5}{L_1^5}\right)\biggr].
 \label{eq:T-log-expansion}
\end{align}
Every finite truncation of \cref{eq:T-log-expansion} is still much larger than
the first outer correction \(a_1/T\).  Indeed,
\begin{equation}
 \frac{a_1/T}{T}=O\!\left(\frac{L_1}{Y}\right)
 =o(L_1^{-M})\qquad(M\ge0).
 \label{eq:outer-beyond-slow}
\end{equation}

\subsection{The fully expanded nested transseries}

Combining \cref{eq:main-transseries,eq:W-log-expansion} gives a conventional
iterated-logarithm transseries of the shape
\begin{equation}
 K_+^{-1}(y)
 \sim \frac12+
 \sum_{n\ge0}2^{1-2n}
 \left(\frac{L_1}{Y}\right)^{n-1/2}
 \sum_{r\ge0}\frac{P_{n,r}(L_2)}{L_1^r},
 \label{eq:nested-log-transseries}
\end{equation}
where each \(P_{n,r}\) is a polynomial.  The convention for \(n=0\) is
\((L_1/Y)^{-1/2}=\sqrt{Y/L_1}\), and \(P_{0,0}=1\).  The first row of these
polynomials is displayed explicitly in \cref{eq:T-log-expansion}.  All later
rows follow mechanically by substituting \cref{eq:W-log-expansion} into
\(a_n((w+1)/2)T^{1-2n}\).

\begin{warningbox}
The double sum in \cref{eq:nested-log-transseries} must be truncated according
to the transseries ordering, not by a naive rectangle.  For every fixed
\(r\), all terms in the leading \(n=0\) row dominate the entire \(n=1\) row.
The \(W\)-anchored form automatically respects this hierarchy and is therefore
the recommended computational representation.
\end{warningbox}

\subsection{Reversing the hierarchy: \(1/\lambda\) first}

Equation \cref{eq:s-implicit} permits the alternative expansion
\begin{equation}
 u(q,\lambda)
 \sim U_0(q)+\frac{U_1(q)}{\lambda}
 +\frac{U_2(q)}{\lambda^2}+\cdots.
 \label{eq:Uj-resummation}
\end{equation}
This resums infinitely many outer blocks at each fixed slow order.

Set
\begin{equation}
 r=r(q):=\sqrt{1+\frac q{12}}.
 \label{eq:r-def}
\end{equation}
Then the leading function is exactly
\begin{equation}
 U_0(q)=r-1.
 \label{eq:U0}
\end{equation}
The next one is also explicit.

\begin{proposition}[First slow correction]
\label{prop:U1}
With the source series interpreted formally in \(q\),
\begin{equation}
 U_1(q)=-\frac1r\left\{
 \frac12\log r+
 \left(12C-\frac14\right)(r^2-1)
 +\sum_{m\ge1}\kappa_mq^{m+1}r^{-2m}
 \right\}.
 \label{eq:U1}
\end{equation}
\end{proposition}

\begin{proof}
Insert \(u=U_0+\lambda^{-1}U_1+\cdots\) into
\cref{eq:s-implicit}.  At order \(\lambda^0\),
\(A(U_0)=q/24\), giving \cref{eq:U0}.  At order \(\lambda^{-1}\),
\begin{equation}
 A'(U_0)U_1+R(U_0,q)=0.
\end{equation}
Since \(A'(U_0)=1+U_0=r\), it remains to simplify \(R(r-1,q)\).  Using
\(q=12(r^2-1)\), one finds
\begin{equation}
 B(r-1)+q\left(C-\frac{\log r}{24}\right)
 =\frac12\log r+\left(12C-\frac14\right)(r^2-1),
\end{equation}
which proves the formula.
\end{proof}

All higher \(U_j\) satisfy a triangular functional recurrence.  For \(j\ge1\),
let
\begin{equation}
 V_{j-1}(s):=U_0+\sum_{i=1}^{j-1}s^iU_i.
 \label{eq:Vj}
\end{equation}
Then
\begin{equation}
 rU_j
 +\frac12\sum_{i=1}^{j-1}U_iU_{j-i}
 +[s^{j-1}]R(V_{j-1}(s),q)=0.
 \label{eq:Uj-recurrence}
\end{equation}
This representation is useful when many outer powers of \(q\) are needed but
only a few inverse powers of \(\lambda\) are desired.

\begin{remark}[Two valid but different summation directions]
The expansion \(u=\sum_n a_n(\lambda)q^n\) is outer-first and is asymptotically
canonical because \(q\ll\lambda^{-M}\).  The expansion
\(u=\sum_j\lambda^{-j}U_j(q)\) is a formal resummation across infinitely many
outer blocks.  It is algebraically useful, but analytic claims about the
\(q\)-series \(U_j\) require separate convergence or summability analysis.
\end{remark}

\section{An accelerated Lambert anchor}
\label{sec:accelerated-anchor}

The leading Lambert anchor treats the constant and logarithmic terms as a
perturbation.  A different elementary model absorbs both of them exactly and
also captures a nontrivial infinite string of inverse powers.  This gives a
markedly sharper closed approximation before any transseries correction is
added.

Set
\begin{equation}
 d:=\frac1{12},
 \qquad
 \widehat\Phi(t)
 :=\frac14(t^2-d)\bigl(\log(t^2-d)-1\bigr)+C.
 \label{eq:Phihat-def}
\end{equation}
For positive large \(t\), expansion of the logarithm gives
\begin{equation}
 \widehat\Phi(t)
 =\Phi(t)+C-\frac1{24}\log t
  +\sum_{m\ge1}\widehat\kappa_m t^{-2m},
 \qquad
 \widehat\kappa_m:=\frac{d^{m+1}}{4m(m+1)}.
 \label{eq:Phihat-expansion}
\end{equation}
Indeed,
\begin{equation}
 (1-x)\log(1-x)
 =-x+\sum_{j\ge2}\frac{x^j}{j(j-1)},
 \label{eq:one-minus-x-log}
\end{equation}
with \(x=d/t^2\), and the constant terms cancel.  Therefore the exact
centered logarithm may be written
\begin{equation}
 \Lcal(t)
 \sim \widehat\Phi(t)+\widehat H(t),
 \qquad
 \widehat H(t):=\sum_{m\ge1}\rho_m t^{-2m},
 \label{eq:accelerated-splitting}
\end{equation}
where
\begin{equation}
 \rho_m:=\kappa_m-\widehat\kappa_m.
 \label{eq:rho-def}
\end{equation}
The first residual coefficients are
\begin{equation}
 \rho_1=-\frac1{480},
 \qquad
 \rho_2=\frac{61}{362880},
 \qquad
 \rho_3=-\frac{433}{4354560}.
 \label{eq:rho-first}
\end{equation}

\subsection{Exact inversion of the accelerated model}

Let \(Y=\log y\), and assume \(Y>C\).  Define
\begin{equation}
 \widehat w:=\W_0\!\left(\frac{4(Y-C)}{\e}\right),
 \qquad
 S:=\frac{4(Y-C)}{\widehat w},
 \qquad
 \widehat T:=\sqrt{S+d}.
 \label{eq:accelerated-anchor-def}
\end{equation}
Then
\begin{equation}
 \widehat\Phi(\widehat T)=Y.
 \label{eq:Phihat-inversion}
\end{equation}
To see this, set \(X=\widehat T^2-d=S\).  The defining equation becomes
\begin{equation}
 X(\log X-1)=4(Y-C),
 \label{eq:X-logX-accelerated}
\end{equation}
and the substitution \(\widehat w=\log X-1\) reduces it to
\(\widehat w\e^{\widehat w}=4(Y-C)/\e\).

\begin{mainbox}[title={Accelerated one-line inverse}]
For \(y\to+\infty\), with \(Y=\log y\),
\begin{equation}
 \boxed{
 K_+^{-1}(y)
 =\frac12+
 \sqrt{\frac{4(Y-C)}{\W_0(4(Y-C)/\e)}+\frac1{12}}
 +O\!\left(\frac{1}{T^3\log T}\right).}
 \label{eq:accelerated-one-line}
\end{equation}
Here \(C=\log A-\tfrac18=-\zeta'(-1)-\tfrac1{24}\), and either \(T\) or
\(\widehat T\) may be used inside the displayed error term.
\end{mainbox}

The gain over the uncorrected anchor \(z\approx\tfrac12+T\) is two outer
powers: the latter has error \(O(T^{-1})\), whereas
\cref{eq:accelerated-one-line} has error \(O(T^{-3}/\log T)\).  The
acceleration is not accidental: it resums the complete \(\lambda^0\) row
\(U_0(q)=\sqrt{1+q/12}-1\), shifts the constant \(C\) into the Lambert
argument, and incorporates the logarithmic correction into the modified
quadratic variable \(t^2-d\).

\subsection{All-orders corrections around the accelerated anchor}

Since
\begin{equation}
 \widehat\Phi'(t)
 =\frac t2\log(t^2-d),
 \label{eq:Phihat-prime}
\end{equation}
define the anchor derivative operator
\begin{equation}
 \widehat{\mathscr D}
 :=\frac{2}{t\log(t^2-d)}\frac{\dd}{\dd t}.
 \label{eq:Dhat-def}
\end{equation}
The perturbative formula of \cref{thm:perturbative-LB} gives the formal
accelerated inverse
\begin{equation}
 t
 \sim \widehat T+
 \sum_{n\ge1}\frac{(-1)^n}{n!}
 \widehat{\mathscr D}^{\,n-1}
 \left.
 \left(
   \frac{\widehat H(t)^n}{\widehat\Phi'(t)}
 \right)\right|_{t=\widehat T}.
 \label{eq:accelerated-LB}
\end{equation}
Every fixed inverse order receives contributions from only finitely many
perturbation orders.  The first correction follows solely from
\(\rho_1=-1/480\):
\begin{equation}
 t
 =\widehat T+
 \frac{1}{240\widehat T^3\log(\widehat T^2-1/12)}
 +O\!\left(\frac{1}{\widehat T^5\log\widehat T}
           +\frac{1}{\widehat T^7(\log\widehat T)^3}\right).
 \label{eq:accelerated-first-correction}
\end{equation}
The two terms displayed in the remainder come respectively from the next
source coefficient and from nonlinear reversion of the first one.

\begin{proposition}[Accelerated error transfer]
\label{prop:accelerated-error}
For \(N\ge0\), let \(\widehat t_N\) denote the truncation of
\cref{eq:accelerated-LB} selected so that substitution into the accelerated
forward transseries cancels every residual term through source order
\(t^{-2N}\); thus \(\widehat t_0=\widehat T\).  Then, on the positive real
axis,
\begin{equation}
 K_+^{-1}(\e^Y)-\frac12-\widehat t_N
 =O_N\!\left(\frac{\widehat T^{-2N-3}}{\log\widehat T}\right),
 \qquad Y\to+\infty.
 \label{eq:accelerated-error}
\end{equation}
For \(N=0\), this is \cref{eq:accelerated-one-line}.  For \(N=1\),
retaining the correction in \cref{eq:accelerated-first-correction} cancels
\(\rho_1t^{-2}\) and moves the leading error to the scale stated inside that
formula's remainder.
\end{proposition}

\begin{proof}
By construction, the accelerated forward residual is
\(O(t^{-2N-2})\).  Division by
\(\widehat\Phi'(t)\asymp t\log t\), followed by the same mean-value argument
as in \cref{thm:real-asymptotic}, gives
\(O(t^{-2N-3}/\log t)\).  The retained formal reversion terms cancel all
larger nonlinear contributions, and \(t\sim\widehat T\) converts this to
\cref{eq:accelerated-error}.
\end{proof}

\begin{remark}[Which anchor should be used?]
The basic anchor is structurally simpler and exposes the polynomial block
coefficients \(a_n(\lambda)\).  The accelerated anchor is usually superior
for a one-line numerical approximation.  They encode the same inverse germ;
they merely choose different portions of the forward transseries as the
exactly invertible model.
\end{remark}

\section{Complex sectorial inverse germs}
\label{sec:complex-branches}

The positive real inverse uses the principal logarithm and \(\W_0\), but the
same construction describes complex inverse germs.  Because both the
logarithm of the target and Lambert \(W\) are multivalued, two branch indices
appear naturally.

Fix a branch of \(\log K(t+\tfrac12)\) in a sector
\begin{equation}
 |\arg t|\le \pi-\delta,
 \qquad \delta>0,
 \label{eq:t-sector}
\end{equation}
where the centered Stirling expansion is uniform.  For a nonzero target
\(y\), choose
\begin{equation}
 Y_\ell:=\Log y+2\pi i\ell,
 \qquad \ell\in\mathbb Z,
 \label{eq:Yell}
\end{equation}
and a branch \(\W_k\) of Lambert \(W\).  Set
\begin{equation}
 w_{k,\ell}:=\W_k\!\left(\frac{4Y_\ell}{\e}\right),
 \qquad
 \lambda_{k,\ell}:=\frac{w_{k,\ell}+1}{2},
 \qquad
 T_{k,\ell}:=\exp(\lambda_{k,\ell}).
 \label{eq:complex-anchor}
\end{equation}
Then \(\Phi(T_{k,\ell})=Y_\ell\) on the corresponding logarithmic sheet.

\begin{theorem}[Sectorial complex inverse transseries]
\label{thm:complex-inverse}
Let a pair \((k,\ell)\) vary along a path such that
\begin{equation}
 |T_{k,\ell}|\longrightarrow\infty,
 \qquad
 |\arg T_{k,\ell}|\le\pi-\delta,
 \qquad
 |w_{k,\ell}+1|\ge\varepsilon>0.
 \label{eq:complex-nondegenerate}
\end{equation}
Assume also that the chosen local solution avoids critical points of
\(\log K\).  Then there is a sectorial inverse germ with
\begin{equation}
 z_{k,\ell}(y)
 \sim \frac12+T_{k,\ell}
 +\sum_{n\ge1}a_n(\lambda_{k,\ell})
 T_{k,\ell}^{1-2n},
 \label{eq:complex-inverse-series}
\end{equation}
where the same rational functions \(a_n\) as on the real branch occur.  For
every fixed \(N\), the error is
\begin{equation}
 O_{N,\delta,\varepsilon}
 \bigl(|T_{k,\ell}|^{-2N-1}\bigr)
 \label{eq:complex-error}
\end{equation}
through any closed subsector on which the derivative remains uniformly
nonzero.
\end{theorem}

\begin{proof}
The centered forward expansion and its differentiated form are uniform in
closed subsectors of \cref{eq:t-sector}; this follows from the corresponding
sectorial Stirling expansions for \(\Gamma\) and Barnes \(G\)
\cite[\S5.11, \S5.17]{DLMF}.  The formal coefficient calculation is
algebraic and is unchanged after complexification.  Condition
\(|w+1|\ge\varepsilon\) excludes the critical value of the Lambert anchor,
for which \(\Phi'(T)=T(w+1)/2\) vanishes.  The analytic inverse-function
theorem and the complex mean-value replacement based on contour estimates
then transfer the residual order to the inverse order.
\end{proof}

\begin{warningbox}
The indices \((k,\ell)\) label asymptotic constructions, not necessarily
globally distinct single-valued branches of an entire inverse.  Analytic
continuation can permute germs, two index pairs can describe the same germ on
overlapping sheets, and singularities occur at critical values of \(K\) as
well as at the Lambert branch point \(w=-1\).  Formula
\cref{eq:complex-inverse-series} is therefore a sectorial local statement.
\end{warningbox}

The accelerated anchor also complexifies.  Replacing \(Y\) by \(Y_\ell\),
choosing a branch \(\W_k(4(Y_\ell-C)/\e)\), and choosing the compatible square
root gives
\begin{equation}
 \widehat T_{k,\ell}^2
 =\frac{4(Y_\ell-C)}{\W_k(4(Y_\ell-C)/\e)}+\frac1{12}.
 \label{eq:complex-accelerated}
\end{equation}
The operator formula \cref{eq:accelerated-LB} then applies on any domain where
\(\log(\widehat T^2-1/12)\) and the square root are consistently chosen and
the anchor derivative does not vanish.

\section{A general theory for Lambert-invertible logarithmic transseries}
\label{sec:general-theory}

The calculation above is an instance of a broad and stable inversion
mechanism.  Its essential feature is a dominant monomial multiplied by one
logarithm.  Lambert \(W\) inverts that dominant phase exactly, while a
differential operator converts every lower-order source monomial into a
strictly smaller inverse monomial.

\subsection{The exact generalized Lambert anchor}

Let
\begin{equation}
 F_0(x):=a x^p(\log x-b),
 \qquad a\ne0,\quad p>0,
 \label{eq:general-F0}
\end{equation}
and solve \(F_0(X)=Y\).  Put
\begin{equation}
 w:=\W_k\!\left(\frac{p}{a}\e^{-pb}Y\right),
 \qquad
 L:=b+\frac wp,
 \qquad
 X:=\e^L
      =\left(\frac{pY}{aw}\right)^{1/p}.
 \label{eq:general-anchor}
\end{equation}
Then \(F_0(X)=Y\).  The logarithmic slope variable is
\begin{equation}
 \Lambda:=p(L-b)+1=w+1,
 \label{eq:Lambda-def}
\end{equation}
and
\begin{equation}
 F_0'(X)=aX^{p-1}\Lambda.
 \label{eq:general-F0-prime}
\end{equation}
Thus the only intrinsic degeneracy of the anchor occurs at the Lambert
critical point \(w=-1\).

\subsection{A one-grid source class}

Fix \(\sigma>0\) and define the outer monomial
\begin{equation}
 Q:=X^{-\sigma}.
 \label{eq:general-Q}
\end{equation}
Consider a source transseries of the form
\begin{equation}
 F(x)=F_0(x)+H(x),
 \qquad
 H(x)=x^{p-\sigma}\Hcal(x^{-\sigma},\log x),
 \label{eq:general-source}
\end{equation}
where
\begin{equation}
 \Hcal(Q,L)=\sum_{m\ge0}h_m(L)Q^m.
 \label{eq:general-Hcal}
\end{equation}
The coefficients \(h_m\) belong to a differential algebra \(\Acal\) that is
closed under \(\dd/\dd L\), multiplication, and division by \(a\Lambda\).
Typical choices are rational functions in \(L\), Laurent polynomials in
\(\Lambda\), or logarithmic transseries in slower scales.

\begin{definition}[Lambert-grid inverse transseries]
\label{def:Lambert-grid}
A formal inverse associated with \cref{eq:general-source} is an expansion
\begin{equation}
 x
 =X\left(1+\sum_{N\ge1}c_N(L)Q^N\right),
 \qquad c_N\in\Acal[\Lambda^{-1}],
 \label{eq:general-inverse-shape}
\end{equation}
ordered first by the powers of \(Q\) and then by the chosen asymptotic order
inside the coefficient algebra.
\end{definition}

The anchor derivative operator is
\begin{equation}
 \mathscr D_0
 :=\frac1{F_0'(X)}\frac{\dd}{\dd X}.
 \label{eq:general-D0}
\end{equation}
For every exponent \(\rho\) and coefficient function \(f(L)\),
\begin{equation}
 \mathscr D_0\bigl(X^\rho f(L)\bigr)
 =X^{\rho-p}\Bcal_\rho f(L),
 \qquad
 \Bcal_\rho f:=\frac{f'(L)+\rho f(L)}{a\Lambda}.
 \label{eq:general-block-action}
\end{equation}
Each application lowers the power of \(X\) by exactly \(p\), while remaining
inside the coefficient differential algebra.

\begin{theorem}[Universal inverse formula]
\label{thm:universal-inverse}
For the source \cref{eq:general-source}, the formal inverse satisfying
\(F(x)=Y=F_0(X)\) is
\begin{equation}
 \boxed{
 x
 \sim X+\sum_{n\ge1}\frac{(-1)^n}{n!}
 \mathscr D_0^{\,n-1}
 \left(\frac{H(X)^n}{F_0'(X)}\right).}
 \label{eq:universal-inverse}
\end{equation}
At each prescribed power of \(Q\), only finitely many values of \(n\)
contribute.  Consequently \cref{eq:universal-inverse} is well defined as a
formal transseries independently of convergence of the perturbation series in
an auxiliary parameter.
\end{theorem}

\begin{proof}
Formula \cref{eq:universal-inverse} is
\cref{eq:perturbative-LB} with the auxiliary parameter set to one.  A monomial
from \(H^n/F_0'\) has initial size
\begin{equation}
 X^{n(p-\sigma)-(p-1)-m\sigma}
 =X^{(n-1)p-n\sigma+1-m\sigma},
 \end{equation}
where \(m\ge0\) is its total source-grid index.  After \(n-1\) applications
of \(\mathscr D_0\), \cref{eq:general-block-action} changes this to
\begin{equation}
 X^{1-(m+n)\sigma}=XQ^{m+n}.
 \label{eq:local-finiteness-power}
\end{equation}
Hence the coefficient of \(XQ^N\) can only receive contributions from
\(1\le n\le N\), with \(m=N-n\).  This proves local finiteness and the claimed
shape.
\end{proof}

\begin{theorem}[Ordered coefficient formula]
\label{thm:general-block-formula}
Let
\begin{equation}
 \rho_{n,N}:=(n-1)p-N\sigma+1.
 \label{eq:rho-nN}
\end{equation}
Then the coefficient in \cref{eq:general-inverse-shape} is
\begin{align}
 c_N(L)
 ={}&\sum_{n=1}^{N}\frac{(-1)^n}{n!}
 \bigl(
   \Bcal_{\rho_{n,N}-(n-2)p}
   \circ\cdots\circ
   \Bcal_{\rho_{n,N}-p}
   \circ\Bcal_{\rho_{n,N}}
 \bigr)
 \notag\\[-0.2em]
 &\hspace{31mm}\times
 \left(
  \frac{[Q^{N-n}]\Hcal(Q,L)^n}{a\Lambda}
 \right).
 \label{eq:general-block-formula}
\end{align}
For \(n=1\), the operator composition is empty.  For \(n\ge2\), it contains
exactly \(n-1\) factors, and the rightmost factor acts first.
\end{theorem}

\begin{proof}
The coefficient \([Q^{N-n}]\Hcal^n\) in \(H^n/F_0'\) is multiplied by
\(X^{\rho_{n,N}}\).  Successive applications of \(\mathscr D_0\) lower its
power by \(p\), so \cref{eq:general-block-action} produces the displayed
ordered list of \(\Bcal\)-operators.  The final power is \(X^{1-N\sigma}\),
as required.
\end{proof}

\subsection{Normalized implicit recurrence}

An alternative construction avoids the perturbative operator formula.  Set
\begin{equation}
 x=X(1+u),
 \qquad
 u=\sum_{N\ge1}c_NQ^N.
 \label{eq:general-u}
\end{equation}
After division by \(aX^p\), the equation \(F(x)=F_0(X)\) becomes
\begin{align}
 0={}&(1+u)^p
 \left(L+\log(1+u)-b\right)-(L-b)
 \notag\\
 &+\frac{Q}{a}(1+u)^{p-\sigma}
 \Hcal\bigl(Q(1+u)^{-\sigma},L+\log(1+u)\bigr).
 \label{eq:general-normalized-implicit}
\end{align}
The derivative of its first line with respect to \(u\) at \(u=0\) is
\(\Lambda\).  Therefore the coefficient of \(c_N\) in the \(Q^N\)-equation
is \(\Lambda\), and the recurrence is triangular:
\begin{equation}
 c_N=-\frac1\Lambda[Q^N]
 \Fcal\!\left(\sum_{j=1}^{N-1}c_jQ^j,Q;L\right),
 \label{eq:general-triangular}
\end{equation}
where \(\Fcal\) denotes the left-hand side of
\cref{eq:general-normalized-implicit}.  Formula
\cref{eq:general-triangular} is often the simplest route for computer
algebra; \cref{eq:general-block-formula} is usually the clearest theoretical
formula.

\subsection{Asymptotic validity}

\begin{theorem}[Residual-to-inverse transfer for the general class]
\label{thm:general-analytic}
Suppose \(F\) is analytic in a real ray or closed complex sector, the source
expansion is differentiable term by term there, and
\begin{equation}
 H(x)=x^{p-\sigma}\sum_{m=0}^{N-1}h_m(\log x)x^{-m\sigma}
      +O\bigl(x^{p-(N+1)\sigma}\Acal_N(x)\bigr),
 \label{eq:general-source-remainder}
\end{equation}
where \(\Acal_N\) is a controlled slow factor.  Assume
\begin{equation}
 |\Lambda|\ge\varepsilon>0,
 \qquad
 |F'(x)|\asymp |a\Lambda X^{p-1}|
 \label{eq:general-nondegenerate}
\end{equation}
uniformly near the desired root.  If \(x_N\) is the inverse transseries
truncated after \(c_NXQ^N\), then
\begin{equation}
 x-x_N
 =O\!\left(X^{1-(N+1)\sigma}\widetilde\Acal_N(X)\right),
 \label{eq:general-inverse-remainder}
\end{equation}
where \(\widetilde\Acal_N\) differs from \(\Acal_N\) only by factors generated
inside the coefficient differential algebra, in particular by powers of
\(\Lambda^{-1}\).
\end{theorem}

\begin{proof}
Substitution of the recursively determined coefficients cancels the normalized
residual through \(Q^N\).  Restoring the factor \(aX^p\) leaves a residual of
order \(X^{p-(N+1)\sigma}\) times a slow factor.  Division by the local slope
\(F'(x)\asymp a\Lambda X^{p-1}\), using either the real mean-value theorem or
a sectorial inverse-function estimate, produces
\(X^{1-(N+1)\sigma}\) and the indicated slow factors.
\end{proof}

\subsection{Beyond a single grid}

The one-grid hypothesis is convenient, not essential.  Suppose instead that
\begin{equation}
 H(x)=\sum_{\gamma\in\Gamma}
 x^{p-\gamma}h_\gamma(\log x),
 \label{eq:multigrid-source}
\end{equation}
where \(\Gamma\subset(0,\infty)\) is well ordered and locally finite.  Products
of source monomials generate the additive semigroup
\begin{equation}
 \langle\Gamma\rangle
 :=\left\{\gamma_1+\cdots+\gamma_n:
 n\ge1,\ \gamma_j\in\Gamma\right\}.
 \label{eq:additive-support}
\end{equation}
The same power count as in \cref{eq:local-finiteness-power} shows that the
relative inverse support is contained in \(\langle\Gamma\rangle\).  Whenever
that semigroup is locally finite---as it is for every finitely generated
positive grid---the perturbative formula is coefficientwise finite.  Hahn or
logarithmic transseries coefficient fields may be used for the slow
coefficients.  This is the natural setting for mixtures of integer powers,
fractional powers, iterated logarithms, and exponentially small sectors; see
\cite{vanDerHoeven2006,ADH2017} for the broader transseries framework.

\begin{methodbox}[title={General Lambert-grid inversion recipe}]
Given a large equation \(F(x)=Y\) whose dominant part is
\(a x^p(\log x-b)\):
\begin{enumerate}
 \item invert the dominant phase exactly by \cref{eq:general-anchor};
 \item choose an outer support for the lower-order source and verify local
       finiteness under addition;
 \item compute coefficients either by the normalized recurrence
       \cref{eq:general-triangular} or by the ordered operator formula
       \cref{eq:general-block-formula};
 \item retain Lambert \(W\) in the final expression unless a fully logarithmic
       expansion is explicitly needed;
 \item certify a finite approximation from its residual and a lower bound on
       \(|F'|\).
\end{enumerate}
\end{methodbox}

\section{Examples and neighboring inverse problems}
\label{sec:examples}

The parameters in the general theory clarify both what is special about the
\(K\)-function and what is universal.  The following table records the dominant
phase, the first outer spacing, and the exact Lambert anchor in several
standard problems.  In each row \(Y\) denotes the logarithm of the original
large target when the special function itself, rather than its logarithm, is
being inverted.

\begin{table}[htbp]
\centering
\small
\caption{Lambert anchors for related logarithmic inverse problems.}
\label{tab:neighboring-anchors}
\begin{tabularx}{\textwidth}{@{}>{\raggedright\arraybackslash}p{0.22\textwidth}
  >{\centering\arraybackslash}p{0.18\textwidth}
  >{\centering\arraybackslash}p{0.12\textwidth}
  >{\centering\arraybackslash}X@{}}
\toprule
Problem & \((a,p,b)\) & \(\sigma\) & Dominant anchor \(X\) \\
\midrule
\(\log\Gamma(x)=Y\)
 & \((1,1,1)\) & \(1\)
 & \(Y/\W(Y/\e)\) \\
Centered \(\log K(t+\tfrac12)=Y\)
 & \((\tfrac12,2,\tfrac12)\) & \(2\)
 & \(2\sqrt{Y/\W(4Y/\e)}\) \\
\(\log G(x+1)=Y\)
 & \((\tfrac12,2,\tfrac32)\) & \(1\)
 & \(2\sqrt{Y/\W(4Y/\e^3)}\) \\
Order-\(r\) hyperfactorial, \(p=r+1\)
 & \((1/p,p,1/p)\) & typically \(1\)
 & \(\bigl(p^2Y/\W(p^2Y/\e)\bigr)^{1/p}\) \\
\bottomrule
\end{tabularx}
\end{table}

\subsection{The inverse gamma function}

Stirling's expansion has the form
\begin{equation}
 \log\Gamma(x)
 \sim x(\log x-1)-\frac12\log x+\frac12\log(2\pi)
 +\sum_{m\ge1}\frac{B_{2m}}{2m(2m-1)x^{2m-1}}.
 \label{eq:gamma-general-example}
\end{equation}
Thus \(a=1\), \(p=1\), \(b=1\), and the correction begins at
\(x^0\), so \(\sigma=1\).  With
\begin{equation}
 X=\frac{Y}{\W(Y/\e)},
 \qquad
 \Lambda=\W(Y/\e)+1,
 \label{eq:gamma-anchor-example}
\end{equation}
the first general coefficient is
\begin{equation}
 c_1=-\frac{h_0(L)}{\Lambda},
 \qquad
 h_0(L)=-\frac L2+\frac12\log(2\pi).
 \label{eq:gamma-c1-example}
\end{equation}
This recovers the familiar Lambert description of the large inverse gamma
branch while placing all subsequent corrections in a single automatic
calculus.

\subsection{The centered inverse \(K\)-function}

For \(x=t\), the centered source has
\begin{equation}
 a=\frac12,
 \qquad p=2,
 \qquad b=\frac12,
 \qquad \sigma=2,
 \qquad h_0(L)=C-\frac L{24}.
 \label{eq:K-general-parameters}
\end{equation}
Since \(a\Lambda=\tfrac12(w+1)=\lambda\), the universal first coefficient
\begin{equation}
 c_1=-\frac{h_0(L)}{a\Lambda}
 \label{eq:universal-c1}
\end{equation}
becomes
\begin{equation}
 c_1=\frac1{24}-\frac C\lambda=a_1(\lambda),
 \label{eq:K-c1-recovered}
\end{equation}
exactly as in \cref{eq:a1}.  The unusually large spacing \(\sigma=2\) is the
benefit of the canonical half-shift and parity cancellation.

\subsection{The inverse Barnes \(G\)-function}

The standard Barnes expansion is
\begin{align}
 \log G(x+1)
 \sim{}&\frac{x^2}{2}\log x-\frac{3x^2}{4}
 +\frac x2\log(2\pi)-\frac1{12}\log x+\zeta'(-1)
 \notag\\
 &+\sum_{m\ge1}\frac{B_{2m+2}}
 {4m(m+1)x^{2m}}.
 \label{eq:BarnesG-example}
\end{align}
The dominant phase has \((a,p,b)=(\tfrac12,2,\tfrac32)\), but the
\(x\log(2\pi)/2\) correction begins one power below the dominant term.
Consequently \(\sigma=1\), and the inverse contains every integer outer power
rather than only the even lattice obtained for the centered \(K\)-function.
This comparison explains why the identity
\begin{equation}
 K(z)G(z)=\Gamma(z)^{z-1}
 \label{eq:K-G-comparison}
\end{equation}
does not by itself reveal the sparse inverse support of \(K\): the sparsity
appears only after cancellation and centering.

\subsection{Higher generalized hyperfactorials}

For an integer \(r\ge0\), consider the discrete products
\begin{equation}
 H_r(n):=\prod_{k=1}^{n}k^{k^r}.
 \label{eq:Hr-discrete}
\end{equation}
Euler--Maclaurin summation gives, with \(p=r+1\),
\begin{equation}
 \log H_r(n)
 =\frac{n^p}{p}\log n-\frac{n^p}{p^2}
  +O(n^{p-1}\log n).
 \label{eq:Hr-leading}
\end{equation}
Any normalized continuous interpolation having this asymptotic form therefore
has the dominant parameters
\begin{equation}
 a=\frac1p,
 \qquad b=\frac1p,
 \qquad
 X=\left(
 \frac{p^2Y}{\W(p^2Y/\e)}
 \right)^{1/p}.
 \label{eq:Hr-anchor}
\end{equation}
The lower Euler--Maclaurin terms normally begin at relative order \(X^{-1}\),
so the generic grid spacing is \(\sigma=1\).  Special shifts can eliminate
some sublattices, just as the half-shift does for \(K\).  Determining such a
shift amounts to choosing a center at which the relevant Bernoulli
polynomials vanish.

\subsection{The first two universal coefficients}

For later reuse, write the source as in \cref{eq:general-Hcal}.  The first two
relative inverse coefficients obtained from
\cref{eq:general-block-formula} are
\begin{align}
 c_1={}&-\frac{h_0}{a\Lambda},
 \label{eq:general-c1}\\
 c_2={}&-\frac{h_1}{a\Lambda}
 +\frac12\Bcal_{p-2\sigma+1}
 \left(\frac{h_0^2}{a\Lambda}\right).
 \label{eq:general-c2}
\end{align}
These formulas already include differentiation of all logarithmic coefficient
functions through \(\Bcal_\rho\).  They are therefore valid without assuming
that the \(h_m\) are constants.

\section{Symbolic generation and stable numerical evaluation}
\label{sec:algorithms}

The recurrence and operator formulas are straightforward to automate.  This
section gives reproducible pseudocode in Python/SymPy notation and a stable
workflow for evaluating the inverse without forming the astronomically large
number \(y=\e^Y\).

\subsection{Generating the centered source coefficients}

The Bernoulli-polynomial formula \cref{eq:kappa-formula} is directly
implementable:

\begin{lstlisting}
import sympy as sp

m = sp.symbols("m", integer=True, positive=True)

def kappa(j: int) -> sp.Expr:
    if j < 1:
        raise ValueError("j must be positive")
    return sp.factor(
        -sp.bernoulli(2*j + 2, sp.Rational(1, 2))
        / (4*j*(j + 1)*(2*j + 1))
    )

print([kappa(j) for j in range(1, 6)])
# [-7/5760, 31/161280, -127/1290240,
#  511/4866048, -1414477/7380172800]
\end{lstlisting}

\subsection{Triangular generation of inverse blocks}

The following routine implements \cref{eq:triangular-recurrence}.  Every
expression is truncated before coefficient extraction, preventing irrelevant
high powers from expanding.

\begin{lstlisting}
import sympy as sp

q, lam, C = sp.symbols("q lam C")

def trunc(expr: sp.Expr, order: int) -> sp.Expr:
    return sp.series(expr, q, 0, order + 1).removeO().expand()

def inverse_blocks(N: int) -> list[sp.Expr]:
    if N < 1:
        return []

    u = sp.Integer(0)
    blocks: list[sp.Expr] = []

    for n in range(1, N + 1):
        log1pu = trunc(sp.log(1 + u), n)
        psi = sp.Rational(1, 2)*(1 + u)**2 * (
            lam + log1pu - sp.Rational(1, 2)
        ) - sp.Rational(1, 2)*(lam - sp.Rational(1, 2))

        source = q*(C - (lam + log1pu)/24)
        for j in range(1, n):
            source += kappa(j)*q**(j + 1) * trunc(
                (1 + u)**(-2*j), n
            )

        residual_n = trunc(psi + source, n).coeff(q, n)
        a_n = sp.factor(-residual_n/lam)
        blocks.append(a_n)
        u = trunc(u + a_n*q**n, n)

    return blocks
\end{lstlisting}

Because the coefficient equation is triangular, no nonlinear equation solver
is invoked.  Exact rational arithmetic is preserved, and \(C\) may remain a
symbol until numerical evaluation.

\subsection{Evaluation from \(Y=\log y\)}

For very large targets, passing \(y\) itself to a floating-point routine is
unnecessary and may overflow.  The natural input is \(Y\).  The basic
Lambert-anchored truncation is evaluated as follows.

\begin{lstlisting}
from mpmath import mp

mp.dps = 80
Cnum = -mp.diff(mp.zeta, -1) - mp.mpf(1)/24

def inverse_K_from_logY(Y, coeff_functions):
    Y = mp.mpf(Y)
    w = mp.lambertw(4*Y/mp.e)
    lam = (w + 1)/2
    T = mp.e**lam

    z = mp.mpf(1)/2 + T
    for n, a_n in enumerate(coeff_functions, start=1):
        z += a_n(lam, Cnum) * T**(1 - 2*n)
    return z
\end{lstlisting}

A robust production implementation should also:
\begin{enumerate}
 \item evaluate \(T=2\sqrt{Y/w}\) rather than \(\e^\lambda\) when that is
       numerically better scaled;
 \item increase working precision before subtracting a residual near zero;
 \item use the exact derivative \cref{eq:logK-prime} for a Newton correction;
 \item certify the final answer with \cref{eq:residual-certificate} whenever a
       rigorous enclosure is needed.
\end{enumerate}

\begin{algorithm}[Lambert-transseries inverse of \(K\)]
\label{alg:inverse-K}
Given \(Y=\log y\), a block order \(N\), and a working precision:
\begin{enumerate}
 \item Compute \(w=\W_0(4Y/\e)\), \(\lambda=(w+1)/2\), and
       \(T=2\sqrt{Y/w}\).
 \item Form
       \(z_N=\tfrac12+T+\sum_{n=1}^{N}a_n(\lambda)T^{1-2n}\).
 \item Optionally replace \(z_N\) by one Newton iterate
       \begin{equation*}
       z_N-\frac{(z_N-1)\log\Gamma(z_N)-\log G(z_N)-Y}
       {z_N-\tfrac12-\tfrac12\log(2\pi)+\log\Gamma(z_N)}.
       \end{equation*}
 \item Estimate or bound the remaining error from the new residual.
\end{enumerate}
\end{algorithm}

The numerator in the Newton step uses the exact positive-real identity
\cref{eq:K-Barnes-relation}; it avoids numerical integration in the defining formula
for \(K\).

\subsection{Choosing a truncation order}

For fixed \(N\), the theorem gives an unambiguous asymptotic order.  If many
blocks are generated, however, the factorial growth of the Bernoulli source
coefficients eventually propagates into the inverse.  A practical adaptive
rule is:
\begin{equation}
 \text{stop before the first block whose magnitude exceeds its predecessor,}
 \label{eq:least-block-rule}
\end{equation}
then apply one Newton step.  This is a numerical heuristic rather than a
replacement for a remainder theorem.  When only moderate precision is needed,
three or four blocks are already far beyond the accuracy of the leading
Lambert approximation over a wide range of \(Y\).

\section{High-precision numerical checks}
\label{sec:numerics}

The formulas were checked against the exact identity
\begin{equation}
 \log K(z)=(z-1)\log\Gamma(z)-\log G(z),
 \qquad z>0,
 \label{eq:logK-exact-numeric}
\end{equation}
using high-precision arithmetic.  The reference root was computed directly
from \(\log K(z)=Y\), and the reported numbers are absolute errors in \(z\).
No Newton correction was applied to the asymptotic values in the tables.

\begin{table}[htbp]
\centering
\small
\caption{Absolute error of the basic Lambert-anchored truncations
\(z_N=\tfrac12+T+\sum_{n=1}^{N}a_nT^{1-2n}\).}
\label{tab:basic-errors}
\begin{tabular}{@{}rrrrrr@{}}
\toprule
\(Y\) & exact inverse \(z_*\) & \(N=0\) & \(N=1\) & \(N=2\) & \(N=4\) \\
\midrule
\(10\)
 & \(4.966157945496709\)
 & \(9.162\times10^{-3}\) & \(1.935\times10^{-5}\)
 & \(1.347\times10^{-7}\) & \(6.963\times10^{-11}\) \\
\(100\)
 & \(10.915136726993085\)
 & \(1.070\times10^{-3}\) & \(2.046\times10^{-7}\)
 & \(5.165\times10^{-10}\) & \(2.775\times10^{-14}\) \\
\(1000\)
 & \(27.284589040600087\)
 & \(1.503\times10^{-4}\) & \(2.134\times10^{-8}\)
 & \(4.241\times10^{-12}\) & \(4.537\times10^{-18}\) \\
\(10^6\)
 & \(584.235864411692401\)
 & \(3.809\times10^{-5}\) & \(2.524\times10^{-13}\)
 & \(6.128\times10^{-19}\) & \(2.420\times10^{-30}\) \\
\bottomrule
\end{tabular}
\end{table}

The \(N=3\) errors, omitted from the table for width, are respectively
\begin{equation}
 6.036\times10^{-10},\qquad
 2.690\times10^{-12},\qquad
 3.032\times10^{-15},\qquad
 8.319\times10^{-25}.
 \label{eq:N3-errors}
\end{equation}
The nonmonotone comparison between \(N=3\) and \(N=4\) at \(Y=10\) is a
normal finite-argument effect; asymptotic truncations are not required to
improve monotonically at every point.

\begin{table}[htbp]
\centering
\small
\caption{Absolute error of the accelerated anchor and of its first correction.}
\label{tab:accelerated-errors}
\begin{tabular}{@{}rrr@{}}
\toprule
\(Y\) & anchor \cref{eq:accelerated-one-line}
& anchor plus \cref{eq:accelerated-first-correction} \\
\midrule
\(10\)    & \(1.559\times10^{-5}\)  & \(6.167\times10^{-8}\) \\
\(100\)   & \(7.865\times10^{-7}\)  & \(5.825\times10^{-10}\) \\
\(1000\)  & \(3.297\times10^{-8}\)  & \(3.706\times10^{-12}\) \\
\(10^6\)  & \(1.644\times10^{-12}\) & \(3.894\times10^{-19}\) \\
\bottomrule
\end{tabular}
\end{table}

Several features are visible.
\begin{itemize}
 \item The first basic correction removes roughly two to eight decimal orders
       over the displayed range, with the gain increasing rapidly with \(Y\).
 \item Four basic blocks deliver about \(14\) correct decimal places already at
       \(Y=100\), and about \(30\) digits at \(Y=10^6\).
 \item The accelerated one-line anchor has nearly the accuracy of the first
       basic correction, while its first residual correction is competitive
       with two or three basic blocks.
\end{itemize}
These numerical observations agree with the proved orders; they are not used
as evidence for the derivation.

\section{Conclusions}
\label{sec:conclusion}

The inverse generalized hyperfactorial is governed by a sparse two-scale
transseries once the correct variable is chosen.  The central facts are:
\begin{enumerate}
 \item The shift \(t=z-\tfrac12\) converts the logarithm of \(K\) into
       \(\Phi(t)+C-\tfrac1{24}\log t\) plus only even inverse powers.
 \item The dominant phase \(\Phi(t)=\tfrac12t^2\log t-\tfrac14t^2\) is
       inverted exactly by \(w=\W(4Y/\e)\), not merely approximated by an
       iterated-logarithm series.
 \item The inverse has the all-orders form
       \(\tfrac12+T+\sum_{n\ge1}a_n(\lambda)T^{1-2n}\), with a triangular
       recurrence, an ordered operator formula, and a fixed-order remainder
       \(O(T^{-2N-1})\).
 \item Every \(K\)-coefficient block is a finite polynomial in
       \(\lambda^{-1}\); the constant row resums to
       \(\sqrt{1+q/12}-1\).
 \item A modified quadratic Lambert model absorbs the constant, logarithm,
       and an infinite rational tail, yielding a substantially sharper
       one-line approximation.
 \item The same construction works for
       \(a x^p(\log x-b)\) plus locally finite lower-order logarithmic
       transseries.  Lambert inversion supplies the anchor, and
       Lagrange--B\"urmann supplies a universal coefficient calculus.
\end{enumerate}

The practical recommendation is equally simple: retain Lambert \(W\) as an
exact structural component, use a small number of outer correction blocks, and
finish with a residual-controlled Newton step when numerical accuracy is the
goal.  For formal work, the normalized recurrence is easiest to automate; for
proofs and support analysis, the ordered differential-operator formula is the
more informative representation.

\appendix

\section{Formula sheet}
\label{app:formula-sheet}

This appendix collects the formulas needed for direct use without repeating
the derivations.

\subsection{Forward data}

\begin{align}
 t&=z-\frac12,
 &\Phi(t)&=\frac12t^2\log t-\frac14t^2,
 &C&=-\zeta'(-1)-\frac1{24}.
 \label{eq:sheet-forward-1}
\end{align}
\begin{equation}
 \log K\!\left(t+\frac12\right)
 \sim \Phi(t)+C-\frac1{24}\log t+
 \sum_{m\ge1}\kappa_m t^{-2m}.
 \label{eq:sheet-forward-2}
\end{equation}
\begin{equation}
 \kappa_m
 =\frac{(1-2^{-2m-1})B_{2m+2}}
 {4m(m+1)(2m+1)}.
 \label{eq:sheet-kappa}
\end{equation}

\subsection{Basic Lambert anchor and inverse}

\begin{align}
 Y&=\log y,
 &w&=\W_0(4Y/\e),
 &\lambda&=\frac{w+1}{2},
 \label{eq:sheet-anchor-1}\\
 T&=2\sqrt{Y/w}=\e^\lambda,
 &q&=T^{-2}=\frac{w}{4Y},
 &&
 \label{eq:sheet-anchor-2}\\
 K_+^{-1}(y)
 &\sim\frac12+T+\sum_{n\ge1}a_n(\lambda)T^{1-2n}.
 &&
 \label{eq:sheet-inverse}
\end{align}
The first block is
\begin{equation}
 a_1=\frac1{24}-\frac C\lambda.
 \label{eq:sheet-a1}
\end{equation}
The higher blocks through \(a_4\) are given in
\cref{eq:a2,eq:a3,eq:a4}.  The normalized recurrence is
\begin{equation}
 a_N=-\frac1\lambda[q^N]\Fcal(U_{N-1},q;\lambda),
 \qquad
 U_{N-1}=\sum_{j=1}^{N-1}a_jq^j,
 \label{eq:sheet-recurrence}
\end{equation}
with \(\Fcal\) defined by \cref{eq:normalized-implicit}.

\subsection{Accelerated anchor}

\begin{align}
 \widehat w&=\W_0\!\left(\frac{4(Y-C)}{\e}\right),
 &\widehat T&=
 \sqrt{\frac{4(Y-C)}{\widehat w}+\frac1{12}},
 \label{eq:sheet-accelerated-1}\\
 K_+^{-1}(y)
 &=\frac12+\widehat T
 +\frac{1}{240\widehat T^3
 \log(\widehat T^2-1/12)}+\cdots.
 \label{eq:sheet-accelerated-2}
\end{align}

\subsection{General Lambert-grid inversion}

For
\begin{equation}
 F(x)=a x^p(\log x-b)
 +x^{p-\sigma}\Hcal(x^{-\sigma},\log x),
 \label{eq:sheet-general-source}
\end{equation}
put
\begin{equation}
 w=\W_k\!\left(\frac pa\e^{-pb}Y\right),
 \qquad
 X=\left(\frac{pY}{aw}\right)^{1/p},
 \qquad
 \Lambda=w+1.
 \label{eq:sheet-general-anchor}
\end{equation}
Then
\begin{equation}
 x\sim X+\sum_{n\ge1}\frac{(-1)^n}{n!}
 \left(\frac1{aX^{p-1}\Lambda}\frac\dd{\dd X}\right)^{n-1}
 \left(
 \frac{H(X)^n}{aX^{p-1}\Lambda}
 \right).
 \label{eq:sheet-general-inverse}
\end{equation}

\section{Conventions and shift dictionary}
\label{app:conventions}

Several neighboring conventions coexist in the literature.  The formulas in
this article use
\begin{equation}
 K(x+1)=x^xK(x),
 \qquad K(1)=1,
 \qquad K(n+1)=\prod_{k=1}^n k^k.
 \label{eq:appendix-K-convention}
\end{equation}
If the hyperfactorial is written
\begin{equation}
 H(x):=K(x+1),
 \label{eq:appendix-H-convention}
\end{equation}
then
\begin{equation}
 H_+^{-1}(y)=K_+^{-1}(y)-1.
 \label{eq:appendix-H-inverse}
\end{equation}
If a discrete source defines \(H_n=\prod_{k=1}^{n-1}k^k\), then
\(H_n=K(n)\), whereas the convention
\(\prod_{k=1}^{n}k^k\) equals \(K(n+1)\).  These shifts do not change the
Lambert scale, but they do change the constant term of an inverse expansion
and therefore must be fixed before comparing formulas.

The Bernoulli convention is \(B_1=-1/2\), generated by
\begin{equation}
 \frac{u}{\e^u-1}=\sum_{n\ge0}B_n\frac{u^n}{n!}.
 \label{eq:appendix-Bernoulli}
\end{equation}
The symbol \(\W_k\) denotes the \(k\)-th branch of Lambert \(W\), characterized
locally by \(\W_k(z)\e^{\W_k(z)}=z\).  On the positive real axis all main
formulas use \(\W_0\).

\section{A complete high-precision verification script}
\label{app:verification-code}

The following compact script regenerates the coefficients by the triangular
recurrence, computes the exact positive-real benchmark through Barnes \(G\),
and reproduces tables of inverse errors.  It uses SymPy for exact symbolic
algebra and mpmath for numerical evaluation.

\begin{lstlisting}
import sympy as sp
from mpmath import mp

q, lam, Csym = sp.symbols("q lam C")

def kap(j):
    return sp.factor(
        -sp.bernoulli(2*j + 2, sp.Rational(1, 2))
        / (4*j*(j + 1)*(2*j + 1))
    )

def trunc(expr, n):
    return sp.series(expr, q, 0, n + 1).removeO().expand()

def make_blocks(N):
    u = sp.Integer(0)
    ans = []
    for n in range(1, N + 1):
        log1pu = trunc(sp.log(1 + u), n)
        psi = sp.Rational(1, 2)*(1 + u)**2 * (
            lam + log1pu - sp.Rational(1, 2)
        ) - sp.Rational(1, 2)*(lam - sp.Rational(1, 2))
        F = psi + q*(Csym - (lam + log1pu)/24)
        for j in range(1, n):
            F += kap(j)*q**(j + 1)*trunc((1 + u)**(-2*j), n)
        an = sp.factor(-trunc(F, n).coeff(q, n)/lam)
        ans.append(an)
        u = trunc(u + an*q**n, n)
    return ans

mp.dps = 100
Cnum = -mp.diff(mp.zeta, -1) - mp.mpf(1)/24
blocks = make_blocks(4)
block_num = [sp.lambdify((lam, Csym), a, "mpmath") for a in blocks]

def logK(z):
    return (z - 1)*mp.loggamma(z) - mp.log(mp.barnesg(z))

def exact_inverse(Y):
    w = mp.lambertw(4*Y/mp.e)
    guess = mp.mpf(1)/2 + 2*mp.sqrt(Y/w)
    return mp.findroot(lambda z: logK(z) - Y, guess)

def basic_inverse(Y, N):
    w = mp.lambertw(4*Y/mp.e)
    la = (w + 1)/2
    T = 2*mp.sqrt(Y/w)
    z = mp.mpf(1)/2 + T
    for n in range(N):
        z += block_num[n](la, Cnum)*T**(-2*n - 1)
    return z

def accelerated_inverse(Y, corrected=False):
    wh = mp.lambertw(4*(Y - Cnum)/mp.e)
    Th = mp.sqrt(4*(Y - Cnum)/wh + mp.mpf(1)/12)
    z = mp.mpf(1)/2 + Th
    if corrected:
        z += 1/(240*Th**3*mp.log(Th**2 - mp.mpf(1)/12))
    return z

for Y in map(mp.mpf, [10, 100, 1000, 10**6]):
    root = exact_inverse(Y)
    print("Y =", Y, "root =", mp.nstr(root, 40))
    print("basic errors:",
          [mp.nstr(abs(basic_inverse(Y, N) - root), 12)
           for N in range(5)])
    print("accelerated errors:",
          [mp.nstr(abs(accelerated_inverse(Y, flag) - root), 12)
           for flag in (False, True)])
\end{lstlisting}

At positive real \(z\), \(G(z)>0\) and the logarithms in the script are
unambiguous.  For complex tests, one must maintain compatible branches of
\(\log\Gamma\), \(\log G\), \(\log K\), and \(\W\) along the continuation
path.

\section{Project context and reproducibility}
\label{app:project-context}

This note was prepared for the ``Series and transseries'' research-frontier
folder of the ProveIt project:
\begin{center}
\url{https://github.com/VladimirReshetnikov/ProveIt/tree/main/Analysis/FabiusFunction/docs/semi-formalized-research-frontiers/drafts/series-and-transseries}.
\end{center}
The \LaTeX{} source is standalone: it has no external figures, bibliography
files, or input fragments.  A standard \texttt{pdflatex} run performed twice
is sufficient to resolve the table of contents and cross-references.

\begin{thebibliography}{99}

\bibitem{ADH2017}
M.~Aschenbrenner, L.~van den Dries, and J.~van der Hoeven,
\emph{Asymptotic Differential Algebra and Model Theory of Transseries},
Annals of Mathematics Studies, vol.~195, Princeton University Press,
Princeton, 2017.
\newblock \url{https://doi.org/10.1515/9781400885411}.

\bibitem{Bendersky1933}
L.~Bendersky,
``Sur la fonction gamma g\'en\'eralis\'ee,''
\emph{Acta Mathematica} \textbf{61} (1933), 263--322.
\newblock \url{https://doi.org/10.1007/BF02547794}.

\bibitem{BrassescoMendez2011}
S.~Brassesco and M.~A. M\'endez,
``The asymptotic expansion for \(n!\) and the Lagrange inversion formula,''
\emph{The Ramanujan Journal} \textbf{24} (2011), 219--234.
\newblock \url{https://doi.org/10.1007/s11139-010-9237-2}.

\bibitem{CorlessEtAl1996}
R.~M. Corless, G.~H. Gonnet, D.~E. Hare, D.~J. Jeffrey, and D.~E. Knuth,
``On the Lambert \(W\) function,''
\emph{Advances in Computational Mathematics} \textbf{5} (1996), 329--359.
\newblock \url{https://doi.org/10.1007/BF02124750}.

\bibitem{DLMF}
NIST Digital Library of Mathematical Functions,
F.~W.~J. Olver, A.~B. Olde Daalhuis, D.~W. Lozier, B.~I. Schneider,
R.~F. Boisvert, C.~W. Clark, B.~R. Miller, B.~V. Saunders,
H.~S. Cohl, and M.~A. McClain, eds.,
\S5.11 (Gamma-function asymptotics), \S5.13 (integrals),
\S5.17 (Barnes \(G\)-function), and \S4.13 (Lambert \(W\)).
\newblock \url{https://dlmf.nist.gov/}.

\bibitem{MathSEHyperfactorial}
C.~Leibovici,
``Looking for the inverse of the hyperfactorial function,''
Mathematics Stack Exchange, question 3507137, 13 January 2020.
\newblock \url{https://math.stackexchange.com/questions/3507137/}.

\bibitem{SondowHadjicostas2007}
J.~Sondow and P.~Hadjicostas,
``The generalized-Euler-constant function \(\gamma(z)\) and a generalization
of Somos's quadratic recurrence constant,''
\emph{Journal of Mathematical Analysis and Applications}
\textbf{332} (2007), 292--314.
\newblock \url{https://doi.org/10.1016/j.jmaa.2006.09.081}.

\bibitem{vanDerHoeven2006}
J.~van der Hoeven,
\emph{Transseries and Real Differential Algebra},
Lecture Notes in Mathematics, vol.~1888, Springer, Berlin, 2006.
\newblock \url{https://doi.org/10.1007/3-540-35590-1}.

\bibitem{Wang2017}
W.~Wang,
``Some asymptotic expansions on hyperfactorial functions and generalized
Glaisher--Kinkelin constants,''
\emph{The Ramanujan Journal} \textbf{43} (2017), 513--533.
\newblock \url{https://doi.org/10.1007/s11139-017-9909-2}.

\bibitem{WangLiu2016}
W.~Wang and H.~Liu,
``Asymptotic expansions related to hyperfactorial function and
Glaisher--Kinkelin constant,''
\emph{Applied Mathematics and Computation} \textbf{283} (2016), 153--162.
\newblock \url{https://doi.org/10.1016/j.amc.2016.02.027}.

\bibitem{Xu2020}
J.~Xu,
``New asymptotic expansions on hyperfactorial functions,''
\emph{Comptes Rendus Math\'ematique} \textbf{358} (2020), no.~9--10,
971--980.
\newblock \url{https://doi.org/10.5802/crmath.73}.

\end{thebibliography}

\end{document}
